\chapter{群论中的进一步主题}

\section{$p$-群、幂零群与可解群}

设 $p$ 是素数，$G$ 是阶为 $p^a n$ 的有限群，其中 $p\nmid n$。回忆，一个（有限）$p$-群是指阶为 $p$ 的幂的群。Sylow 定理说明，$p$-群大量出现在 $G$ 的子群中；为了利用这一现象来剖析有限群的结构，我们需要先建立 $p$-群的一些基本性质。下一节中我们将在许多具体情形中应用这些结果。

在给出关于 $p$-群的结果之前，先回忆一个在前面练习中已经出现过的定义。

\medskip
\noindent\textbf{定义.} 群 $G$ 的一个\emph{极大子群}是指 $G$ 的一个真子群 $M$，使得不存在 $G$ 的子群 $H$ 满足
\[
  M < H < G.
\]

由阶数考虑可知，有限群的每个真子群都包含在某个极大子群中。相反，无限群可以有极大子群，也可以没有极大子群。例如，$p\ZZ$ 是 $\ZZ$ 的极大子群，而 $\QQ$（加法群）没有极大子群（参见本节末练习 16）。

我们现在把以后需要的 $p$-群性质集中到一个总括性定理中。

\medskip
\noindent\textbf{定理 1.} 设 $p$ 是素数，$P$ 是阶为 $p^a$ 的群，其中 $a\geq 1$。则：
\begin{enumerate}[label=(\arabic*)]
  \item $P$ 的中心非平凡：$Z(P)\neq 1$。
  \item 若 $H$ 是 $P$ 的非平凡正规子群，则 $H$ 与中心的交非平凡：$H\cap Z(P)\neq 1$。特别地，每个阶为 $p$ 的正规子群都包含在中心中。
  \item 若 $H\normal P$，则对 $|H|$ 的每个因子 $p^b$，$H$ 含有一个阶为 $p^b$ 且在 $P$ 中正规的子群。特别地，对每个 $b\in\{0,1,\ldots,a\}$，$P$ 都有一个阶为 $p^b$ 的正规子群。
  \item 若 $H<P$，则 $H<N_P(H)$，即 $P$ 的每个真子群都是其在 $P$ 中正规化子的真子群。
  \item $P$ 的每个极大子群都具有指标 $p$，并且在 $P$ 中正规。
\end{enumerate}

\begin{proof}
这些结果归根结底依赖类方程，读者可先回顾第 4.3 节。

(1) 是第 4 章定理 8；它也是 (2) 在 $H=P$ 时的特例。我们先证明 (2)，尽管下面的论证只是第 4 章论证的轻微推广，但这里不直接引用第 4 章定理 8。设 $H$ 是 $P$ 的非平凡正规子群。回忆：对 $P$ 的每个共轭类 $C$，由于 $H$ 正规，或者 $C\subseteq H$，或者 $C\cap H=\varnothing$。取 $P$ 的共轭类代表元
\[
  a_1,a_2,\ldots,a_r,
\]
其中 $a_1,\ldots,a_k\in H$，而 $a_{k+1},\ldots,a_r\notin H$。令 $C_i$ 为 $a_i$ 在 $P$ 中的共轭类。于是
\[
  C_i\subseteq H\quad (1\leq i\leq k),
  \qquad
  C_i\cap H=\varnothing\quad (k+1\leq i\leq r).
\]
必要时重新编号 $a_1,\ldots,a_k$，可假设 $a_1,\ldots,a_s$ 代表大小为 $1$ 的共轭类（即它们在 $Z(P)$ 中），而 $a_{s+1},\ldots,a_k$ 代表大小大于 $1$ 的共轭类。由于 $H$ 是这些类的不交并，
\[
  |H|=|H\cap Z(P)|+\sum_{i=s+1}^{k}|P:C_P(a_i)|.
\]
现在 $p$ 整除 $|H|$，并且 $p$ 整除和式中的每一项 $|P:C_P(a_i)|$，所以 $p$ 整除它们的差 $|H\cap Z(P)|$。因此 $H\cap Z(P)\neq 1$。若 $|H|=p$，则由 $H\cap Z(P)\neq 1$ 得 $H\leq Z(P)$。这证明了 (2)。

接着用关于 $a$ 的归纳法证明 (3)。若 $a=1$ 或 $H=1$，结论平凡。设 $a>1$ 且 $H\neq 1$。由 (2)，$H\cap Z(P)\neq 1$；由 Cauchy 定理，$H\cap Z(P)$ 含有一个阶为 $p$ 的（正规）子群 $Z$。用横线表示取商 $P/Z$。这个商群阶为 $p^{a-1}$，且 $\overline H\normal \overline P$。由归纳假设，对每个使 $p^b$ 整除 $|\overline H|$ 的非负整数 $b$，$\overline H$ 中存在阶为 $p^b$ 且在 $\overline P$ 中正规的子群 $\overline K$。若 $K$ 是 $\overline K$ 在 $P$ 中的完全原像，则 $|K|=p^{b+1}$。由这种过程得到的 $H$ 的所有子群，再加上恒等子群，便给出 $|H|$ 的每个因子阶的、在 $P$ 中正规的子群。(3) 的第二个断言是取 $H=P$ 的特例。

(4) 也用关于 $|P|$ 的归纳法证明。若 $P$ 是阿贝尔群，则 $P$ 的所有子群都在 $P$ 中正规，结论平凡。因此可假设 $|P|>p$（事实上由推论 4.9 可知 $|P|>p^2$）。设 $H$ 是 $P$ 的真子群。由于 $Z(P)$ 的每个元素都与 $P$ 的每个元素交换，所以 $Z(P)$ 正规化 $P$ 的每个子群。由 (1)，$Z(P)\neq 1$。若 $Z(P)\nleq H$，则 $H$ 真包含于 $\langle H,Z(P)\rangle$，而后者又包含在 $N_P(H)$ 中，于是 (4) 成立。故可假设 $Z(P)\leq H$。用横线表示取商 $P/Z(P)$。由于 $\overline P$ 的阶比 $P$ 小，由归纳假设，$\overline H$ 真包含于 $N_{\overline P}(\overline H)$。由格同构定理可知，$N_P(H)$ 是 $N_{\overline P}(\overline H)$ 在 $P$ 中的完全原像，因此在这种情形下也得到 $H$ 真包含于其正规化子。

最后证明 (5)。设 $M$ 是 $P$ 的极大子群。由定义 $M<P$，所以由 (4)，$M<N_P(M)$。由极大性，必有 $N_P(M)=P$，即 $M\normal P$。格同构定理说明 $P/M$ 是一个没有非平凡真子群的 $p$-群；但由 (3)，$P/M$ 对每个整除 $|P/M|$ 的阶都有子群。因此唯一可能是 $|P/M|=p$。证明完毕。
\end{proof}

\medskip
\noindent\textbf{定义.}
\begin{enumerate}[label=(\arabic*)]
  \item 对任意（有限或无限）群 $G$，归纳定义如下子群：
  \[
    Z_0(G)=1,
    \qquad
    Z_1(G)=Z(G),
  \]
  并且 $Z_{i+1}(G)$ 是 $G$ 中包含 $Z_i(G)$ 的子群，使得
  \[
    Z_{i+1}(G)/Z_i(G)=Z(G/Z_i(G)),
  \]
  即 $Z_{i+1}(G)$ 是自然投影下 $G/Z_i(G)$ 的中心在 $G$ 中的完全原像。子群链
  \[
    Z_0(G)\leq Z_1(G)\leq Z_2(G)\leq\cdots
  \]
  称为 $G$ 的\emph{上中心列}。“上”这个词表示 $Z_i(G)\leq Z_{i+1}(G)$。
  \item 若存在 $c\in\ZZ$ 使 $Z_c(G)=G$，则称群 $G$ 是\emph{幂零的}。这样的最小 $c$ 称为 $G$ 的\emph{幂零类}。
\end{enumerate}

本节末的一个练习说明，对所有 $i$，$Z_i(G)$ 都是 $G$ 的特征子群（从而正规）。以后我们将自由使用这一事实。

\paragraph{注}
\begin{enumerate}[label=(\arabic*)]
  \item 若 $G$ 是阿贝尔群，则 $G$ 是幂零的（当 $|G|>1$ 时幂零类为 $1$），因为此时 $G=Z(G)=Z_1(G)$。应把幂零群看作位于阿贝尔群与可解群之间的一类结构：
  \[
    \text{循环群}\subset \text{阿贝尔群}\subset \text{幂零群}\subset \text{可解群}\subset \text{所有群}.
  \]
  上述包含关系均为真包含，稍后将验证。

  \item 对任何有限群，由阶数考虑，必存在整数 $n$ 使得
  \[
    Z_n(G)=Z_{n+1}(G)=Z_{n+2}(G)=\cdots.
  \]
  例如，对所有 $n\in\ZZ^+$，$Z_n(S_3)=1$。一旦上中心列中相邻两项相同，这条链就从该处稳定。例如，若 $G=\ZZ_2\times S_3$，则对所有 $n$，
  \[
    Z(G)=Z_1(G)=Z_2(G)=Z_n(G)
  \]
  的阶都是 $2$。按定义，对非幂零群，所有 $Z_n(G)$ 都是 $G$ 的真子群。

  \item 对无限群 $G$，可能所有 $Z_i(G)$ 都是 $G$ 的真子群（因此 $G$ 不是幂零群），但
  \[
    G=\bigcup_{i=0}^{\infty} Z_i(G).
  \]
  满足这种性质的群称为\emph{超幂零群}；它们具有幂零群的一些性质，但并非全部。虽然我们主要处理有限幂零群，但不涉及阶、Sylow 子群等概念的结果也适用于无限群。即使对无限群，处理幂零群的主要技巧之一也是关于幂零类的归纳法。
\end{enumerate}

\medskip
\noindent\textbf{命题 2.} 设 $p$ 是素数，$P$ 是阶为 $p^a$ 的群。则 $P$ 是幂零群，并且幂零类至多为 $a-1$。

\begin{proof}
对每个 $i\geq 0$，$P/Z_i(P)$ 是 $p$-群。因此，由定理 1(1)，
\[
  |P/Z_i(P)|>1\quad\Longrightarrow\quad Z(P/Z_i(P))\neq 1.
\]
于是若 $Z_i(P)\neq P$，则 $|Z_{i+1}(P)|\geq p|Z_i(P)|$，所以 $|Z_{i+1}(P)|\geq p^{i+1}$。特别地，$|Z_a(P)|\geq p^a$，故 $P=Z_a(P)$，所以 $P$ 的幂零类至多为 $a$。

$P$ 的幂零类恰等于 $a$ 的唯一可能情形是 $|Z_i(P)|=p^i$ 对所有 $i$ 成立。此时 $Z_{a-2}(P)$ 在 $P$ 中指标为 $p^2$，故 $P/Z_{a-2}(P)$ 是阿贝尔群（由推论 4.9）。于是 $P/Z_{a-2}(P)$ 等于其中心，从而 $Z_{a-1}(P)=P$，矛盾。因此 $P$ 的幂零类至多为 $a-1$。
\end{proof}

\paragraph{例} $D_8$ 与 $Q_8$ 都是幂零类为 $2$ 的幂零群。更一般地，$D_{2^n}$ 的幂零类为 $n-1$。这可通过归纳证明：对 $n\geq 3$，$|Z(D_{2^n})|=2$ 且 $D_{2^n}/Z(D_{2^n})\cong D_{2^{n-1}}$（细节留作练习）。若 $n$ 不是 $2$ 的幂，则 $D_{2n}$ 不是幂零群（参见练习 10）。

下面给出有限群幂零性的若干等价（也常常更可操作的）刻画。

\medskip
\noindent\textbf{定理 3.} 设 $G$ 是有限群，$p_1,p_2,\ldots,p_s$ 是整除 $|G|$ 的不同素数，且 $P_i\in\Syl_{p_i}(G)$，$1\leq i\leq s$。则以下条件等价：
\begin{enumerate}[label=(\arabic*)]
  \item $G$ 是幂零群；
  \item 若 $H<G$，则 $H<N_G(H)$，即 $G$ 的每个真子群都是其在 $G$ 中正规化子的真子群；
  \item 对 $1\leq i\leq s$，$P_i\normal G$，即每个 Sylow 子群都在 $G$ 中正规；
  \item $G\cong P_1\times P_2\times\cdots\times P_s$。
\end{enumerate}

\begin{proof}
(1) 推出 (2) 的证明与 $p$-群情形相同；所需的唯一事实是：若 $G$ 幂零，则 $G/Z(G)$ 也幂零。因此细节略去（参见练习）。

为证 (2) 推出 (3)，令 $P=P_i$，$N=N_G(P)$。由于 $P\normal N$，由推论 4.20，$P$ 是 $N$ 的特征子群。因为 $P\charac N\normal N_G(N)$，得到 $P\normal N_G(N)$。这意味着 $N_G(N)\leq N$，从而 $N_G(N)=N$。由 (2)，必有 $N=G$，即 $P\normal G$。

接着证明 (3) 推出 (4)。对任意 $t$，$1\leq t\leq s$，用归纳法证明
\[
  P_1P_2\cdots P_t\cong P_1\times P_2\times\cdots\times P_t.
\]
由于每个 $P_i$ 都在 $G$ 中正规，$P_1\cdots P_t$ 是 $G$ 的子群。令 $H=P_1\cdots P_{t-1}$，$K=P_t$。由归纳假设，$H\cong P_1\times\cdots\times P_{t-1}$，特别地 $|H|=|P_1|\cdots |P_{t-1}|$。而 $|K|=|P_t|$，所以 $|H|$ 与 $|K|$ 互素；由 Lagrange 定理，$H\cap K=1$。又 $P_1\cdots P_t=HK$，由直积识别定理（定理 5.9），
\[
  HK\cong H\times K\cong (P_1\times\cdots\times P_{t-1})\times P_t.
\]
取 $t=s$ 得 (4)。

最后证明 (4) 推出 (1)。由第 5.1 节练习 1，
\[
  Z(P_1\times\cdots\times P_s)\cong Z(P_1)\times\cdots\times Z(P_s).
\]
由第 5.1 节练习 14，
\[
  G/Z(G)\cong (P_1/Z(P_1))\times\cdots\times(P_s/Z(P_s)).
\]
因此 (4) 的假设也适用于 $G/Z(G)$。由定理 1，若 $P_i\neq 1$ 则 $Z(P_i)\neq 1$，所以若 $G\neq 1$，则 $|G/Z(G)|<|G|$。由归纳法，$G/Z(G)$ 幂零；再由练习 6，$G$ 幂零。证明完毕。
\end{proof}

由上一定理立即得到有限阿贝尔群基本定理第一部分（第 5.2 节定理 5）。稍后我们还会由中国剩余定理给出另一个证明。

\medskip
\noindent\textbf{推论 4.} 有限阿贝尔群是它的 Sylow 子群的直积。

下面的命题稍后会用来证明有限域的乘法群是循环群（无需使用有限阿贝尔群基本定理）。

\medskip
\noindent\textbf{命题 5.} 若有限群 $G$ 满足：对每个整除 $|G|$ 的正整数 $n$，方程 $x^n=1$ 在 $G$ 中至多有 $n$ 个解，则 $G$ 是循环群。

\begin{proof}
设 $|G|=p_1^{a_1}\cdots p_s^{a_s}$，并令 $P_i\in\Syl_{p_i}(G)$。由于 $p_i^{a_i}\mid |G|$ 且 $P_i$ 的 $p_i^{a_i}$ 个元素都是 $x^{p_i^{a_i}}=1$ 的解，由假设 $P_i$ 必含有 $G$ 中该方程的所有解。于是 $P_i$ 是唯一的（故正规）Sylow $p_i$-子群。由定理 3，$G$ 是其 Sylow 子群的直积。

由定理 1，每个 $P_i$ 有一个指标为 $p_i$ 的正规子群 $M_i$。由于 $|M_i|=p_i^{a_i-1}$，且 $G$ 中至多有 $p_i^{a_i-1}$ 个元素满足 $x^{p_i^{a_i-1}}=1$，由 Lagrange 定理可知，$M_i$ 含有所有满足该方程的元素。因此 $P_i\setminus M_i$ 中的任一元素都满足 $x^{p_i^{a_i}}=1$ 但 $x^{p_i^{a_i-1}}\neq 1$，也就是说它的阶为 $p_i^{a_i}$。故每个 $P_i$ 都是循环群。$G$ 是两两互素阶循环群的直积，因而是循环群。
\end{proof}

下一个命题称为 Frattini 论证。我们将用它给出有限幂零群的另一个刻画，它在下一节也很有用。

\medskip
\noindent\textbf{命题 6（Frattini 论证）.} 设 $G$ 是有限群，$H\normal G$，$P$ 是 $H$ 的 Sylow $p$-子群。则
\[
  G=HN_G(P),
\]
并且 $|G:H|$ 整除 $|N_G(P)|$。

\begin{proof}
由推论 3.15，$HN_G(P)$ 是 $G$ 的子群，且由于 $H\normal G$，有 $HN_G(P)=N_G(P)H$。设 $g\in G$。由于 $P^g\leq H^g=H$，$P$ 与 $P^g$ 都是 $H$ 的 Sylow $p$-子群。由 Sylow 定理应用于 $H$，存在 $x\in H$ 使 $P^g=P^x$。于是 $gx^{-1}\in N_G(P)$，故 $g\in N_G(P)x$。因为 $g$ 任意，得到 $G=N_G(P)H$。

对 $G=N_G(P)H$ 应用第二同构定理，得到
\[
  |G:H|=|N_G(P):N_G(P)\cap H|,
\]
所以 $|G:H|$ 整除 $|N_G(P)|$。证明完毕。
\end{proof}

\medskip
\noindent\textbf{命题 7.} 有限群幂零当且仅当它的每个极大子群都是正规子群。

\begin{proof}
设 $G$ 是有限幂零群，$M$ 是 $G$ 的极大子群。与定理 1 的证明类似，由定理 3(2) 有 $M<N_G(M)$，而 $M$ 的极大性迫使 $N_G(M)=G$，即 $M\normal G$。

反过来，假设有限群 $G$ 的每个极大子群都是正规子群。令 $P$ 是 $G$ 的一个 Sylow $p$-子群。我们证明 $P\normal G$，然后由定理 3(3) 得 $G$ 幂零。若 $P$ 在 $G$ 中不正规，取一个包含 $N_G(P)$ 的极大子群 $M$。由假设 $M\normal G$，于是由 Frattini 论证，$G=MN_G(P)$。但 $N_G(P)\leq M$，故 $MN_G(P)=M$，矛盾。证明完毕。
\end{proof}

\subsection*{换位子与下中心列}

为完整起见，我们给出群的下中心列的定义，并说明它与上中心列的关系。由于以后不会使用这些结果，证明留作（直接的）练习。

回忆，群 $G$ 中两个元素 $x,y$ 的\emph{换位子}定义为
\[
  [x,y]=x^{-1}y^{-1}xy,
\]
而 $G$ 的两个子群 $H,K$ 的换位子子群为
\[
  [H,K]=\langle [h,k]\mid h\in H,\ k\in K\rangle.
\]
换位子及换位子子群的基本性质已在第 5.4 节建立。

\medskip
\noindent\textbf{定义.} 对任意（有限或无限）群 $G$，归纳定义
\[
  G^0=G,
  \qquad
  G^1=[G,G],
  \qquad
  G^{i+1}=[G,G^i].
\]
子群链
\[
  G=G^0\geq G^1\geq G^2\geq\cdots
\]
称为 $G$ 的\emph{下中心列}。“下”表示 $G^i\geq G^{i+1}$。

与上中心列一样，本节末练习中也要求验证对所有 $i$，$G^i$ 都是 $G$ 的特征子群。下一定理给出上、下中心列之间的关系。

\medskip
\noindent\textbf{定理 8.} 群 $G$ 幂零当且仅当存在 $n\geq 0$ 使 $G^n=1$。更精确地，$G$ 的幂零类为 $c$ 当且仅当 $c$ 是使 $G^c=1$ 的最小非负整数。若 $G$ 的幂零类为 $c$，则对所有 $i\in\{0,1,\ldots,c-1\}$，
\[
  Z_i(G)\leq G^{c-i-1}\leq Z_{i+1}(G).
\]

\begin{proof}
这可通过对上中心列或下中心列的长度作直接归纳证明。
\end{proof}

一般而言，上中心列与下中心列的各项不必相同，不过在某些群中它们会相同。

\paragraph{注}
\begin{enumerate}[label=(\arabic*)]
  \item 若 $G$ 是阿贝尔群，我们已经看到 $G'=G^1=1$，故下中心列一项后终止于单位群。
  \item 与上中心列类似，对任意有限群，由阶数考虑，必存在整数 $n$ 使得
  \[
    G^n=G^{n+1}=G^{n+2}=\cdots.
  \]
  对非幂零群，$G^n$ 是 $G$ 的非平凡子群。例如，在第 5.4 节中我们证明 $S_3'=S_3''=A_3$。由于 $S_3$ 不是幂零的，必有 $S_3^i=A_3$（$i\geq 1$）。事实上
  \[
    (123)=[(12),(132)]\in [S_3,S_3']=S_3^2.
  \]
  一旦下中心列中相邻两项相同，这条链就从该处稳定。注意 $S_3$ 给出一个例子：其下中心列有两个不同项，而上中心列的所有项都等于单位群；因此对非幂零群，两种中心列的长度不必相同。
\end{enumerate}

\subsection*{可解群与导出列}

回忆第 3.4 节中，可解群被定义为具有一个列
\[
  1=H_0\normal H_1\normal\cdots\normal H_s=G,
\]
使每个因子 $H_{i+1}/H_i$ 都是阿贝尔群。现在我们用一列递降的特征子群给出可解性的另一个刻画。

\medskip
\noindent\textbf{定义.} 对任意群 $G$，归纳定义子群列
\[
  G^{(0)}=G,
  \qquad
  G^{(1)}=[G,G],
  \qquad
  G^{(i+1)}=[G^{(i)},G^{(i)}]\quad(i\geq 1).
\]
此子群列称为 $G$ 的\emph{导出列}或\emph{换位子列}。这些项也常写作 $G^{(1)}=G'$、$G^{(2)}=G''$ 等。同样，证明每个 $G^{(i)}$ 都是 $G$ 的特征子群留作练习。

重要的是，虽然 $G^{(0)}=G^0$ 且 $G^{(1)}=G^1$，但一般并不成立 $G^{(i)}=G^i$。区别在于：下中心列的第 $i+1$ 项是第 $i$ 项与整个群 $G$ 的换位子，而导出列的第 $i+1$ 项是第 $i$ 项与自身的换位子。因此
\[
  G^{(i)}\leq G^i\quad\text{对所有 }i,
\]
且包含可以是真包含。例如在 $G=S_3$ 中，$G^1=G'=A_3$ 且 $G^2=[S_3,A_3]=A_3$，而 $G^{(2)}=[A_3,A_3]=1$。

\medskip
\noindent\textbf{定理 9.} 群 $G$ 可解当且仅当存在 $n\geq 0$ 使 $G^{(n)}=1$。

\begin{proof}
先设 $G$ 可解，故存在列
\[
  1=H_0\normal H_1\normal\cdots\normal H_s=G,
\]
且每个因子 $H_{i+1}/H_i$ 都是阿贝尔群。我们用归纳法证明 $G^{(i)}\leq H_{s-i}$。$i=0$ 时显然。若 $G^{(i)}\leq H_{s-i}$，则
\[
  G^{(i+1)}=[G^{(i)},G^{(i)}]\leq [H_{s-i},H_{s-i}].
\]
由于 $H_{s-i}/H_{s-i-1}$ 是阿贝尔群，由命题 5.7(4)，$[H_{s-i},H_{s-i}]\leq H_{s-i-1}$。归纳完成。因 $H_0=1$，得到 $G^{(s)}=1$。

反过来，若 $G^{(n)}=1$，则由命题 5.7(4)，取 $H_i=G^{(n-i)}$，便得到一条每个商群都为阿贝尔的正规列。因此 $G$ 可解。
\end{proof}

若 $G$ 可解，使 $G^{(n)}=1$ 的最小非负整数 $n$ 称为 $G$ 的\emph{可解长度}。导出列是具有阿贝尔逐次商的最短列，并且它还具有所有项都在整个群中为特征子群的性质（而不仅仅是每一项在下一项中正规）。其内在定义也使许多场合更容易处理，如下面命题所示。

\medskip
\noindent\textbf{命题 10.} 设 $G,K$ 是群，$H\leq G$，$\varphi:G\to K$ 是满同态。
\begin{enumerate}[label=(\arabic*)]
  \item 对所有 $i\geq 0$，$H^{(i)}\leq G^{(i)}$。特别地，若 $G$ 可解，则 $H$ 可解；也就是说，可解群的子群可解（且 $H$ 的可解长度不超过 $G$ 的可解长度）。
  \item 对所有 $i\geq 0$，$\varphi(G^{(i)})=K^{(i)}$。特别地，可解群的同态像和商群可解（可解长度不超过定义域的可解长度）。
  \item 若 $N\normal G$，且 $N$ 与 $G/N$ 都可解，则 $G$ 可解。
\end{enumerate}

\begin{proof}
(1) 由 $H\leq G$ 可知 $[H,H]\leq [G,G]$，即 $H^{(1)}\leq G^{(1)}$；再归纳得 $H^{(i)}\leq G^{(i)}$。若 $G^{(n)}=1$，则 $H^{(n)}=1$。

(2) 注意由换位子定义
\[
  \varphi([x,y])=[\varphi(x),\varphi(y)].
\]
故归纳得 $\varphi(G^{(i)})\leq K^{(i)}$。由于 $\varphi$ 是满射，$K$ 中的每个换位子都来自 $G$ 中某个换位子的像，于是归纳得到等号。若 $G^{(n)}=1$，则 $K^{(n)}=1$。

(3) 设 $G/N$ 与 $N$ 的可解长度分别为 $n,m$。对自然投影 $\varphi:G\to G/N$ 应用 (2)，得
\[
  \varphi(G^{(n)})=(G/N)^{(n)}=1N,
\]
即 $G^{(n)}\leq N$。于是
\[
  G^{(n+m)}=(G^{(n)})^{(m)}\leq N^{(m)}=1.
\]
由定理 9，$G$ 可解。证明完毕。
\end{proof}

关于有限群可解性，还有一些附加条件如下。

\medskip
\noindent\textbf{定理 11.} 设 $G$ 是有限群。
\begin{enumerate}[label=(\arabic*)]
  \item （Burnside）若 $|G|=p^aq^b$，其中 $p,q$ 是素数，则 $G$ 可解。
  \item （Philip Hall）若对每个整除 $|G|$ 的素数 $p$，将 $|G|$ 写为 $p^a m$ 且 $(p,m)=1$，并且 $G$ 有阶为 $m$ 的子群，则 $G$ 可解。也就是说，若对每个素数 $p$，$G$ 都有一个指标等于 Sylow $p$-子群阶的子群，则 $G$ 可解；这样的子群称为 Sylow $p$-补。
  \item （Feit--Thompson）若 $|G|$ 是奇数，则 $G$ 可解。
  \item （Thompson）若对每一对元素 $x,y\in G$，$\langle x,y\rangle$ 都是可解群，则 $G$ 可解。
\end{enumerate}

Burnside 定理将在第 19 章证明，并由此推出 Philip Hall 的推广。正如第 3.5 节提到的，Feit--Thompson 定理的证明长达 255 页。Thompson 定理最早是作为一篇 475 页论文的推论得到的，而该论文最终也依赖 Feit--Thompson 定理。

\subsection*{有限阿贝尔群基本定理的一个证明}

我们概述一个群论证明：每个有限阿贝尔群都是循环群的直积（即第 5.2 节定理 5 的第 1、2 部分）。有限生成阿贝尔群分类定理（第 5.2 节定理 3）将在第 12 章作为一个更一般定理的推论得到。

由推论 4，只需证明对素数 $p$，任意阿贝尔 $p$-群都是循环群的直积（第 5.5 定理中的整除条件可通过重排因子平凡地实现）。设 $A$ 是阿贝尔 $p$-群。我们对 $|A|$ 作归纳。

若 $E$ 是初等阿贝尔 $p$-群（即对所有 $x\in E$，$x^p=1$），先证明：
\[
  \text{对任意 }x\in E，\text{存在 }M\leq E\text{ 使 }E=M\times\langle x\rangle.
\]
若 $x=1$，取 $M=E$。否则，取 $M$ 为满足 $x\notin M$ 的、阶尽可能大的 $E$ 的子群。若 $M$ 在 $E$ 中指标不是 $p$，令 $\overline E=E/M$。则 $\overline E$ 是初等阿贝尔群，且存在 $\overline y\in \overline E-\langle \overline x\rangle$。由于 $\overline y$ 阶为 $p$，$\overline x\notin \langle\overline y\rangle$。$\langle\overline y\rangle$ 在 $E$ 中的完全原像是不含 $x$ 且阶大于 $M$ 的子群，矛盾。因此 $|E:M|=p$，于是
\[
  E=M\langle x\rangle,
  \qquad
  M\cap\langle x\rangle=1.
\]
由直积识别定理，$E=M\times\langle x\rangle$。

现在令 $\varphi:A\to A$ 由 $\varphi(x)=x^p$ 定义（见第 5.2 节练习 7）。由于 $A$ 阿贝尔，$\varphi$ 是同态。记 $K=\ker\varphi$，$H=\im\varphi$。则 $K=\{x\in A\mid x^p=1\}$，而 $H$ 是 $A$ 中所有 $p$ 次幂组成的子群。注意 $K$ 与 $A/H$ 都是初等阿贝尔群。由第一同构定理，
\[
  |A:H|=|K|.
\]
由归纳假设，
\[
  H=\langle h_1\rangle\times\cdots\times\langle h_r\rangle
  \cong \ZZ_{p^{a_1}}\times\cdots\times\ZZ_{p^{a_r}},
  \qquad a_i\geq 1.
\]
由 $\varphi$ 的定义，存在 $g_i\in A$ 使 $g_i^p=h_i$。令
\[
  A_0=\langle g_1,\ldots,g_r\rangle.
\]
可验证
\begin{enumerate}[label=(\alph*)]
  \item $A_0=\langle g_1\rangle\times\cdots\times\langle g_r\rangle$；
  \item $A_0/H=\langle g_1H\rangle\times\cdots\times\langle g_rH\rangle$ 是阶为 $p^r$ 的初等阿贝尔群；
  \item $H\cap K=\langle h_1^{p^{a_1-1}}\rangle\times\cdots\times\langle h_r^{p^{a_r-1}}\rangle$ 是阶为 $p^r$ 的初等阿贝尔群。
\end{enumerate}
若 $K\leq H$，则 $|K|=|K\cap H|=p^r=|A_0:H|$，比较阶得到 $A_0=A$，定理得证。

否则 $K\nleq H$。用横线表示取商 $A/H$，取 $x\in K-H$，则 $|\overline x|=|x|=p$。由证明开头对初等阿贝尔 $p$-群的断言，应用于 $\overline E=\overline A$，存在 $\overline A$ 的子群 $\overline M$ 使
\[
  \overline A=\overline M\times\langle\overline x\rangle.
\]
若 $M$ 是 $\overline M$ 在 $A$ 中的完全原像，则由于 $x$ 的阶为 $p$ 且 $x\notin M$，有 $\langle x\rangle\cap M=1$。由直积识别定理，
\[
  A=M\times\langle x\rangle.
\]
由归纳假设，$M$ 是循环群的直积，因此 $A$ 也是循环群的直积。证明完成。

有限阿贝尔群分解为循环群直积的唯一性（第 5.5 定理第 3 部分）也可以用 $p$ 次幂映射作归纳证明。这基本上就是第 12.1 节证明有限生成阿贝尔群基本定理唯一性部分时采用的过程。

\subsection*{练习}
\begin{enumerate}[label=\arabic*.]
  \item 证明对所有 $i$，$Z_i(G)$ 是 $G$ 的特征子群。
  \item 证明定理 1 的 (2) 和 (4) 对任意有限幂零群 $G$ 也成立，不必要求 $G$ 是 $p$-群。
  \item 若 $G$ 有限，证明 $G$ 幂零当且仅当对每个整除 $|G|$ 的阶都有一个正规子群；并证明 $G$ 循环当且仅当对每个整除 $|G|$ 的阶都有唯一子群。
  \item 证明有限幂零群的极大子群具有素数指标。
  \item 证明定理 1 的 (2) 和 (4) 对无限幂零群 $G$ 也成立。
  \item 证明若 $G/Z(G)$ 幂零，则 $G$ 幂零。
  \item 证明幂零群的子群和商群都是幂零的（你的证明应适用于无限群）。给出一个明确例子：某群 $G$ 有正规子群 $H$，且 $H$ 与 $G/H$ 都幂零，但 $G$ 不幂零。
  \item 证明若 $p$ 是素数且 $P$ 是阶为 $p^3$ 的非阿贝尔群，则 $|Z(P)|=p$ 且 $P/Z(P)\cong \ZZ_p\times\ZZ_p$。
  \item 证明有限群 $G$ 幂零当且仅当任取 $a,b\in G$ 且 $(|a|,|b|)=1$ 时都有 $ab=ba$。[使用定理 3 的第 (4) 部分。]
  \item 证明 $D_{2n}$ 幂零当且仅当 $n$ 是 $2$ 的幂。[使用练习 9。]
  \item 在额外假设 $G$ 阿贝尔的情形下，通过调用有限阿贝尔群基本定理给出命题 5 的另一个证明。
  \item 求 $A_4$ 和 $S_4$ 的上中心列与下中心列。
  \item 求 $A_n$ 和 $S_n$（$n\geq 5$）的上中心列与下中心列。
  \item 证明对所有 $i$，$G^i$ 是 $G$ 的特征子群。
  \item 证明 $Z_i(D_{2^n})=D_{2^n}^{\,n-1-i}$。
  \item 证明 $\QQ$ 没有极大子群。[回忆第 3.2 节练习 21。]
  \item 证明对所有 $i$，$G^{(i)}$ 是 $G$ 的特征子群。
  \item 证明若 $G'/G''$ 与 $G''/G'''$ 都是循环群，则 $G''=1$。[可假设 $G'''=1$，然后令 $G/G''$ 通过共轭作用于循环群 $G''$。]
  \item 证明不存在一个群，其换位子子群同构于 $S_4$。[使用前一题。]
  \item 设 $p$ 是素数，$P$ 是有限群 $G$ 的 $p$-子群，$N\normal G$ 且 $|N|$ 与 $p$ 互素，令 $\overline G=G/N$。证明：
  \begin{enumerate}[label=(\alph*)]
    \item $N_{\overline G}(\overline P)=\overline{N_G(P)}$；[使用 Frattini 论证。]
    \item $C_{\overline G}(\overline P)=\overline{C_G(P)}$。[使用 (a)。]
  \end{enumerate}
\end{enumerate}

对任意群 $G$，其 Frattini 子群（记为 $\Phi(G)$）定义为 $G$ 所有极大子群的交；若 $G$ 没有极大子群，则规定 $\Phi(G)=G$。下面几题讨论这个重要子群。

\begin{enumerate}[label=\arabic*.,resume]
  \item 证明 $\Phi(G)$ 是 $G$ 的特征子群。
  \item 证明若 $N\normal G$，则 $\Phi(N)\leq \Phi(G)$。给出一个明确例子说明若 $N$ 不正规，则这个包含不一定成立。
  \item 计算 $\Phi(S_3)$、$\Phi(A_4)$、$\Phi(S_4)$、$\Phi(A_5)$ 和 $\Phi(S_5)$。
  \item 称 $G$ 的元素 $x$ 是\emph{非生成元}，若对 $G$ 的每个真子群 $H$，$\langle x,H\rangle$ 仍是 $G$ 的真子群。证明 $\Phi(G)$ 正是 $G$ 的非生成元集合（这里 $|G|>1$）。
  \item 设 $G$ 是有限群。证明 $\Phi(G)$ 幂零。[用 Frattini 论证证明 $\Phi(G)$ 的每个 Sylow 子群都在 $G$ 中正规。]
  \item 设 $p$ 是素数，$P$ 是有限 $p$-群，令 $\overline P=P/\Phi(P)$。
  \begin{enumerate}[label=(\alph*)]
    \item 证明 $\overline P$ 是初等阿贝尔 $p$-群。[证明 $P'\leq \Phi(P)$ 且对所有 $x\in P$，$x^p\in\Phi(P)$。]
    \item 证明若 $N\normal P$ 且 $P/N$ 是初等阿贝尔群，则 $\Phi(P)\leq N$。用同态与交换图表述这一泛性质。
    \item 设 $\overline P$ 是阶为 $p^r$ 的初等阿贝尔群。由练习 24 推出：若 $\overline x_1,\ldots,\overline x_r$ 是 $\overline P$ 在 $\FF_p$ 上的一个基，且 $x_i$ 是陪集 $\overline x_i$ 中任一元素，则 $P=\langle x_1,
\ldots,x_r\rangle$。反过来证明若 $y_1,\ldots,y_s$ 生成 $P$，则 $s\geq r$。推出 Burnside 基定理：$y_1,
\ldots,y_s$ 是 $P$ 的极小生成集当且仅当其像是 $\overline P=P/\Phi(P)$ 的一个基。并推出 $P$ 的任意极小生成集有 $r$ 个元素。
    \item 证明若 $P/\Phi(P)$ 循环，则 $P$ 循环。推出若 $P/P'$ 循环，则 $P$ 循环。
    \item 设 $\alpha$ 是 $P$ 的素数阶 $q$ 的自同构，且 $q\neq p$。证明若 $\alpha$ 固定陪集 $x\Phi(P)$，则它固定该陪集中的某个元素。[把 $\alpha$ 看作该陪集的 $p^a$ 个元素上的一个阶为 $1$ 或 $q$ 的置换。]
    \item 用 (e) 和 (c) 推出：$P$ 的任意阶与 $p$ 互素的非平凡自同构，在 $P/\Phi(P)$ 上诱导非平凡自同构。推出任意阶与 $p$ 互素的 $P$ 的自同构群同构于 $\Aut(\overline P)=\GL_r(\FF_p)$ 的一个子群。
  \end{enumerate}
  \item 推广前一题 (d)：设 $p$ 是素数，$P$ 是 $p$-群，且 $\overline P=P/\Phi(P)$ 是阶为 $p^r$ 的初等阿贝尔群。证明 $P$ 恰有
  \[
    \frac{p^r-1}{p-1}
  \]
  个极大子群。[由于 $P$ 的每个极大子群都含有 $\Phi(P)$，由格同构定理，$P$ 的极大子群与初等阿贝尔群 $\overline P$ 的极大子群一一对应。因此只需证明阶为 $p^r$ 的初等阿贝尔 $p$-群的极大子群个数如上。一个办法是使用阿贝尔群同构于其对偶群的结果（参见第 5.2 节练习 14），所以指标为 $p$ 的子群个数等于阶为 $p$ 的子群个数。]
  \item 证明若 $p$ 是素数，$P=\ZZ_p\times\ZZ_{p^2}$，则 $|\Phi(P)|=p$ 且 $P/\Phi(P)\cong \ZZ_p\times\ZZ_p$。推出 $P$ 有 $p+1$ 个极大子群。
  \item 证明若 $p$ 是素数，$P$ 是阶为 $p^3$ 的非阿贝尔群，则 $\Phi(P)=Z(P)$ 且 $P/\Phi(P)\cong \ZZ_p\times\ZZ_p$。推出 $P$ 有 $p+1$ 个极大子群。
  \item 令 $p$ 是奇素数，$P_1=\ZZ_p\times\ZZ_{p^2}$，$P_2$ 为有一个阶为 $p^2$ 元素的非阿贝尔 $p^3$ 阶群。证明 $P_1$ 与 $P_2$ 有相同的子群格。
  \item 对任意群 $G$，称 $G$ 的极小正规子群为 $G$ 的一个正规子群 $M$，使得包含在 $M$ 中的 $G$ 的正规子群只有 $1$ 和 $M$。证明有限可解群的每个极小正规子群都是某个素数 $p$ 的初等阿贝尔 $p$-群。[若 $M$ 是 $G$ 的极小正规子群，考察它的特征子群 $M'$ 与 $\langle x^p\mid x\in M\rangle$。]
  \item 证明有限可解群的每个极大子群都有素数幂指标。[令 $H$ 是 $G$ 的极大子群，$M$ 是 $G$ 的极小正规子群——参见前一练习。对 $G/M$ 用归纳法，并分别考虑 $M\leq H$ 与 $M\nleq H$ 两种情形。]
  \item 设 $\pi$ 是任意素数集合。有限群 $G$ 的子群 $H$ 称为 Hall $\pi$-子群，若 $|H|$ 的所有素因子都在 $\pi$ 中，且 $|H|$ 与 $|G:H|$ 互素。（注意若 $\pi=\{p\}$，Hall $\pi$-子群就是 Sylow $p$-子群；Hall 子群在第 3.3 节练习 10 中已引入。）证明可解群的 Sylow 定理推广：若 $G$ 是有限可解群，则对每个素数集合 $\pi$，$G$ 有 Hall $\pi$-子群，且任意两个 Hall $\pi$-子群（对同一集合 $\pi$）共轭。[固定 $\pi$，对 $|G|$ 归纳，并同时证明存在性与共轭性。令 $M$ 是 $G$ 的一个极小正规子群，则 $M$ 是某素数 $p$ 的 $p$-群。若 $p\in\pi$，对 $G/M$ 应用归纳假设。若 $p\notin\pi$，化到 $|G|=p^a n$ 的情形，其中 $p^a=|M|$，$n$ 是 $G$ 的 Hall $\pi$-子群的阶。在这种情况下，令 $N/M$ 是 $G/M$ 的一个极小正规子群，则 $N/M$ 是某个素数 $q\neq p$ 的 $q$-群。令 $Q\in\Syl_q(N)$。若 $Q\normal G$，用 $Q$ 代替 $M$ 继续论证；若 $Q$ 在 $G$ 中不正规，用 Frattini 论证说明 $N_G(Q)$ 是 $G$ 的一个 Hall $\pi$-子群，并在这种情形下也证明共轭性。]
  \item 设 $p$ 整除有限可解群 $G$ 的阶。假设 $G$ 没有非平凡的、阶与 $p$ 互素的正规子群。令 $P$ 是 $G$ 的最大正规 $p$-子群（参见第 4.5 节练习 37）。注意上面练习 31 说明 $P\neq1$。证明 $C_G(P)\leq P$，即 $C_G(P)=Z(P)$。[令 $N=C_G(P)$，并用前一练习证明 $N=Z(P)\times H$，其中 $H$ 是 $N$ 的一个 Hall $\pi$-子群——这里 $\pi$ 是 $|N|$ 的所有不同于 $p$ 的素因子集合。证明 $H\normal G$，从而得到所需结论：$H=1$。]
  \item 证明若有限群 $G$ 的每个真子群都是幂零群，则 $G$ 可解。[证明一个极小反例是单群。取不同极大子群 $M,N$，使 $|M\cap N|$ 尽可能大，并用定理 3 的第 (2) 部分证明 $M\cap N=1$。然后应用第 4.5 节练习 53 的方法。]
  \item 设 $p$ 是素数，$V$ 是 $\FF_p$ 上非零有限维向量空间，$\phi\in\GL(V)$ 的阶为 $p$ 的幂。证明存在非零 $v\in V$ 使 $\phi(v)=v$。
  \item 设 $V$ 是 $\FF_2$ 上有限维向量空间，$\phi\in\GL(V)$ 的阶为 $2$（即 $V$ 是非平凡初等阿贝尔 $2$-群，$\phi$ 是 $V$ 的阶为 $2$ 的自同构）。证明映射 $v\mapsto v+\phi(v)$ 是 $V$ 到自身的同态。证明其像中的每个元素都被 $\phi$ 固定。推出 $\phi$ 的不动点子空间维数至少为 $\frac12\dim V$。（注意若 $G$ 是 $V$ 与 $\langle\phi\rangle$ 的半直积，其中 $V\normal G$ 且 $\phi$ 通过共轭作用在 $V$ 上，把每个 $v\in V$ 送到 $\phi(v)$，则 $\phi$ 在 $V$ 上的不动点就是 $C_V(\phi)$，而上面的映射就是换位子映射 $v\mapsto [v,\phi]$。用这种术语，本题是要证明 $|C_V(\phi)|^2\geq |V|$。）
  \item 用前一题证明：若 $P$ 是中心循环的 $2$-群，$M$ 是 $P$ 中指标为 $2$ 的子群，则 $Z(M)$ 的秩至多为 $2$。[阶为 $2$ 的群 $G/M$ 通过共轭作用在 $\FF_2$ 向量空间 $\{z\in Z(M)\mid z^2=1\}$ 上，而该作用的不动点在 $P$ 的中心中。]
\end{enumerate}

\section{中等阶群中的应用}

本节的目的是通过一系列例子展示前面发展出的许多技巧。这些例子大量使用 Sylow 定理，并说明它们如何用于研究有限群。受 H\"older 纲领启发，我们主要讨论如下问题：对某些 $n$，证明每个阶为 $n$ 的群都有真而非平凡的正规子群（即不存在阶为 $n$ 的单群）。多数情况下，一旦得到这一结论我们就停止。不过读者应注意，在达到这一结果的过程中，我们实际上已经获得了关于给定阶 $n$ 的任意群的大量信息；这些信息可以进一步用于分类这些阶的群，但一般需要超出半直积构造的技术。

由于对素数 $p$，我们已经证明除 $\ZZ_p$ 外不存在单 $p$-群，而 $p$-群的结构又可能非常复杂（回忆第 5.3 节的表），所以我们不显式研究 $p$-群本身的结构。相反，我们将把上一节的 $p$-群理论应用于非素数幂阶群的子群。

最后，对某些 $n$（例如 $60,168,360,504,\ldots$），确实存在阶为 $n$ 的单群，因此当然不能强迫这些阶的所有群都非单。与第 4.5 节一样，我们有时可以证明阶为 $n$ 的单群唯一，并剖析其内部结构（Sylow 数等）。我们将把阶为 $168$ 的单群作为另一个检验例子来研究。因此，Sylow 定理将在不同情境中被应用，以展示如何操纵给定阶的群。

本节末我们还将讨论群的存在性问题，特别是有限单群的存在性问题。

对 $n<10000$，存在 $60$ 个奇的非素数幂整数，使得 Sylow 定理的同余条件不能迫使至少一个 Sylow 子群正规；也就是说，对所有素数 $p\mid n$，$n_p$ 都可能大于 $1$。例如，由第 4.5 节的结果，形如 $pq$（$p,q$ 为不同素数）的数不会出现在我们的名单中。相反，对小于 $500$ 的偶数，已经有 $46$ 个可能成为单群阶的候选数。我们许多数值例子来自这些名单，且常用奇数，因为 Sylow 同余条件给出的 $n_p$ 可能值较少。这些例子的目的在于说明已证明结果的用法；其中许多也可由更高级的技术处理（例如 Feit--Thompson 定理说明不存在奇合数阶单群）。

正如第 4.5 节 $n=30$ 的情形中看到的，尽管 Sylow 定理允许 $n_5=6$ 和 $n_3=10$，进一步分析表明阶为 $30$ 的任意群都必须有 $n_5=1$ 且 $n_3=1$。因此，Sylow 定理的同余部分是 Sylow 子群正规的充分条件，但远非必要条件。对于许多 $n$（如 $n=120$），我们可证明不存在阶为 $n$ 的单群，故存在非平凡正规子群，但该子群未必是 Sylow 子群。例如 $S_5$ 与 $\SL_2(\FF_5)$ 都有阶 $120$；$S_5$ 有唯一的非平凡真正规子群 $A_5$，阶为 $60$，而 $\SL_2(\FF_5)$ 有唯一的非平凡真正规子群 $Z(\SL_2(\FF_5))\cong\ZZ_2$，阶为 $2$；二者都不是 Sylow 子群。因此产生正规子群的技术必须足够灵活。本节中我们将考察不同素数的 Sylow 子群、Sylow 子群的交、$p$-子群的正规化子以及许多不那么显然的子群。这里列出的初等方法即使对“中等”阶群也远非穷尽。

\subsection*{一些技巧}

在列出给定（“中等”）阶群中产生正规子群的一些技巧之前，我们注意：所有处理阶为某个具体 $n$ 的群的问题，首先都必须把 $n$ 分解成素数幂，并计算对每个素数 $p\mid n$，$n_p$ 的允许值。最后一步中要熟练进行模 $p$ 计算。我们描述的技巧如下：
\begin{enumerate}[label=(\arabic*)]
  \item 计数元素；
  \item 利用小指标子群；
  \item 置换表示；
  \item 让不同素数 $p$ 的 $p$-子群相互制约；
  \item 研究 Sylow $p$-子群交的正规化子。
\end{enumerate}

\subsection*{计数元素}

设 $G$ 是阶为 $n$ 的群，$p\mid n$，$P\in\Syl_p(G)$。若 $|P|=p$，则 $P$ 的每个非单位元素阶为 $p$，而 $G$ 中每个阶为 $p$ 的元素都在 $P$ 的某个共轭中。由 Lagrange 定理，不同共轭的 $P$ 只在单位元处相交。因此在此情形下，$G$ 中阶为 $p$ 的元素数为
\[
  n_p(p-1).
\]

若对不同素数 $p$ 的 Sylow $p$-子群都具有素数阶，且假设没有一个正规，则有时可推出素数阶元素个数超过 $|G|$。这种矛盾说明至少一个 $n_p$ 必须为 $1$，即某个 Sylow 子群正规。

这正是第 4.5 节用来证明不存在阶为 $30$ 的单群的论证。再如，设 $|G|=105=3\cdot5\cdot7$。若 $G$ 单，则必须有
\[
  n_3=7,
  \qquad n_5=21,
  \qquad n_7=15.
\]
于是
\[
\begin{array}{rcl}
\text{阶为 }3\text{ 的元素数} &=& 7\cdot2=14,\\
\text{阶为 }5\text{ 的元素数} &=& 21\cdot4=84,\\
\text{阶为 }7\text{ 的元素数} &=& 15\cdot6=90,
\end{array}
\]
总计 $188>|G|$，矛盾。

有时计数素数阶元素不会得到“太多”元素。不过，剩下的元素可能少到必须形成某个正规子群。第 4.5 节证明阶为 $12$ 的群中 $n_2=1$ 或 $n_3=1$ 的方法本质上就是如此。当 $G$ 有一个阶为 $p$ 的 Sylow $p$-子群 $P$ 且 $N_G(P)=P$ 时，这种技巧尤其有效。例如，令 $|G|=56$。若 $G$ 单，则 Sylow $7$-子群数唯一可能是 $8$，所以阶为 $7$ 的元素数为
\[
  8\cdot6=48.
\]
于是 $G$ 中还剩 $56-48=8$ 个元素。Sylow $2$-子群含有 $8$ 个元素，且其中无阶为 $7$ 的元素，因此 Sylow $2$-子群至多一个，故 $G$ 有正规 Sylow $2$-子群。

\subsection*{利用小指标子群}

回忆第 4.2 节的结果：若 $G$ 有指标为 $k$ 的子群 $H$，则存在一个同态 $G\to S_k$，其核包含于 $H$。若 $k>1$，该核是 $G$ 的真正规子群；若我们试图证明 $G$ 非单，可反证假设该核为单位群。于是由第一同构定理，$G$ 同构于 $S_k$ 的一个子群，特别地 $|G|$ 整除 $k!$。因此，若 $G$ 是有限单群，且 $k$ 是满足 $|G|\mid k!$ 的最小整数，则 $G$ 不含指标小于 $k$ 的真子群。研究给定阶 $n$ 的群时，应在一开始就计算这个最小允许指标。

在我们考虑的例子中，这个计算通常很容易：$n$ 常常分解为
\[
  n=p_1^{\alpha_1}p_2^{\alpha_2}\cdots p_s^{\alpha_s},
  \qquad p_1<p_2<\cdots<p_s,
\]
并且在我们的例子中，$\alpha_s$ 通常等于 $1$ 或 $2$。在这种情况下，真子群的最小指标必须至少是 $p_s$（分别地，至少是 $2p_s$），而且这常常就是它的精确值。

例如，不存在阶为 $3393$ 的单群。因为
\[
  3393=3^2\cdot 13\cdot29,
\]
而真子群的最小可能指标是 $29$（$3393\nmid 28!$，因为 $29\nmid 28!$）。但任意阶为 $3393$ 的单群必须有 $n_3=13$，因此对 $P\in\Syl_3(G)$，$N_G(P)$ 的指标为 $13$，矛盾。

\subsection*{置换表示}

这是前一方法的细化。若 $G$ 是阶为 $n$ 的单群，且有指标为 $k$ 的真子群，则 $G$ 同构于 $S_k$ 的一个子群。于是可把 $G$ 看作 $S_k$ 的子群，并在 $S_k$ 内计算。有时不仅仅用 Lagrange 定理，而是直接证明 $S_k$ 不含阶为 $n$ 的单子群。两种有用限制是：
\begin{enumerate}[label=(\arabic*)]
  \item 若 $G$ 含有某个阶的元素或子群，则 $S_k$ 也必须含有；
  \item 若 $P\in\Syl_p(G)$ 且 $P$ 也是 $S_k$ 的 Sylow $p$-子群，则 $|N_G(P)|$ 必须整除 $|N_{S_k}(P)|$。
\end{enumerate}

情形 (2) 常发生在 $p$ 为素数且 $k=p$ 或 $p+1$ 时。此时 $p^2\nmid k!$，所以 $G$ 的 Sylow $p$-子群也是 $S_k$ 的 Sylow $p$-子群。$S_k$ 的 Sylow $p$-子群正是由一个 $p$-轮换生成的群，且不同 Sylow $p$-子群交平凡。因此
\[
  \#\Syl_p(S_k)=
  \frac{\text{$p$-轮换个数}}{\text{一个 Sylow $p$-子群中的 $p$-轮换个数}}
  =\frac{k(k-1)\cdots(k-p+1)}{p(p-1)}.
\]
该数给出 $S_k$ 中 Sylow $p$-子群正规化子的指标。因此，当 $k=p$ 或 $p+1$ 时，
\[
  |N_{S_k}(P)|=p(p-1).
\]

例如，若 $G$ 是阶为 $396=2^2\cdot3^2\cdot11$ 的单群，则必须有 $n_{11}=12$。若 $P\in\Syl_{11}(G)$，则 $|G:N_G(P)|=12$，所以 $|N_G(P)|=33$。由于 $G$ 有指标 $12$ 的子群，$G$ 同构于 $S_{12}$ 的子群。但在 $S_{12}$ 中，$P\in\Syl_{11}(S_{12})$ 且 $|N_{S_{12}}(P)|=110$。由 $N_G(P)\leq N_{S_{12}}(P)$，应有 $33\mid110$，不可能。因此不存在阶为 $396$ 的单群。

有时可稍进一步，在 $A_k$ 中而非 $S_k$ 中工作。这个小改进只偶尔有用，且只针对偶数阶群。

\medskip
\noindent\textbf{命题 12.}
\begin{enumerate}[label=(\arabic*)]
  \item 若 $G$ 没有指标为 $2$ 的子群且 $G\leq S_k$，则 $G\leq A_k$。
  \item 若 $P\in\Syl_p(S_k)$，其中 $p$ 是奇素数，则 $P\in\Syl_p(A_k)$ 且
  \[
    |N_{A_k}(P)|=\frac12|N_{S_k}(P)|.
  \]
\end{enumerate}

\begin{proof}
(1) 由第二同构定理：若 $G\nleq A_k$，则 $A_k<G A_k$，故 $GA_k=S_k$。于是
\[
  2=|S_k:A_k|=|GA_k:A_k|=|G:G\cap A_k|,
\]
这说明 $G$ 有指标为 $2$ 的子群。

(2) 若 $P\in\Syl_p(S_k)$ 且 $p$ 为奇素数，则由 (1)（或阶数考虑）$P\leq A_k$，故 $P\in\Syl_p(A_k)$。由 Frattini 论证，
\[
  S_k=N_{S_k}(P)A_k.
\]
因此 $N_{S_k}(P)\nleq A_k$，从而 $N_{S_k}(P)\cap A_k=N_{A_k}(P)$ 是 $N_{S_k}(P)$ 的指标 $2$ 子群。
\end{proof}

例如，不存在阶为 $264$ 的单群。若 $G$ 是阶为 $264=2^3\cdot3\cdot11$ 的单群，则必须有 $n_{11}=12$。如前，$G$ 同构于 $S_{12}$ 的子群。由于 $G$ 单（故无指标 $2$ 子群），$G\leq A_{12}$。令 $P\in\Syl_{11}(G)$。由 $n_{11}=12=|G:N_G(P)|$，有 $|N_G(P)|=22$。而
\[
  |N_{A_{12}}(P)|=\frac12|N_{S_{12}}(P)|=\frac12\cdot11(11-1)=55.
\]
但 $22\nmid55$，与 $N_G(P)\leq N_{A_{12}}(P)$ 矛盾。

最后，我们强调，我们只是略微触及了某些置换表示所能提供的组合信息。在余下的例子中，只要可能，我们还会说明这一技巧的其他应用。

\subsection*{让不同素数的 $p$-子群相互制约}

设 $p,q$ 是不同素数，且每个阶为 $pq$ 的群都是循环群。若 $p<q$，这等价于 $p\nmid q-1$。若 $G$ 有阶为 $q$ 的 Sylow $q$-子群 $Q$，且 $p\mid |N_G(Q)|$，则在 $N_G(Q)$ 中由 Cauchy 定理得到一个阶为 $p$ 的群 $P$ 正规化 $Q$（注意 $P$ 不一定是 $G$ 的 Sylow $p$-子群）。于是 $PQ$ 是群；若 $PQ$ 阿贝尔，则
\[
  Q\leq N_G(P),
  \qquad \text{从而 } q\mid |N_G(P)|.
\]
类似地，若 $G$ 的 Sylow $p$-子群阶为 $p$，且 $q$ 整除某个 Sylow $p$-正规化子的阶，也可作对称论证。这类数值信息可能足以迫使 $N_G(P)=G$（即 $P\normal G$），或者至少迫使 $N_G(P)$ 具有小于置换表示所允许最小指标的指标，从而产生矛盾。

例如，不存在阶为 $1785$ 的单群。若存在，令 $G$ 是阶为
\[
  1785=3\cdot5\cdot7\cdot17
\]
的单群。$n_{17}$ 的唯一可能值为 $35$，所以若 $Q$ 是 Sylow $17$-子群，则 $|G:N_G(Q)|=35$，故 $|N_G(Q)|=3\cdot17$。令 $P$ 是 $N_G(Q)$ 的 Sylow $3$-子群。群 $PQ$ 是阿贝尔的，因为 $3\nmid 17-1$，所以 $Q\leq N_G(P)$，从而 $17\mid |N_G(P)|$。此时 $P\in\Syl_3(G)$。$n_3$ 的允许值为 $7,85,595$；但由于 $17\mid |N_G(P)|$，不能有 $17\mid |G:N_G(P)|=n_3$，故 $n_3=7$。这意味着 $G$ 有一个指标 $7$ 的真子群，与此阶的最小真子群指标为 $17$ 矛盾。另一种说法是：若 $n_3=7$，则
\[
  |N_G(P)|=3\cdot5\cdot17,
\]
并且由 Sylow 定理应用于 $N_G(P)$，有 $Q\normal N_G(P)$。这与 $|N_G(Q)|=3\cdot17$ 矛盾。

还可以细化这一方法，不要求 $P$ 与 $Q$ 都为素数阶。若 $p,q$ 是整除 $|G|$ 的不同素数，$Q\in\Syl_q(G)$，且 $p\mid |N_G(Q)|$，令 $P\in\Syl_p(N_G(Q))$。可在 $N_G(Q)$ 中应用 Sylow 定理，判断是否 $P\normal N_G(Q)$；若是，则可能迫使 $N_G(P)$ 具有小指标。若 $P$ 是整个 $G$ 的 Sylow $p$-子群，则可再用 Sylow 同余条件限制 $|N_G(P)|$。若 $P$ 不是 $G$ 的 Sylow $p$-子群，则 $P$ 包含在某个 $P^*\in\Syl_p(G)$ 中；由于 $P<P^*$，定理 1(4) 给出 $P<N_{P^*}(P)$，从而 $N_G(P)$ 的阶被比 $|P|$ 更高的 $p$ 的幂整除，并且还被 $|Q|$ 整除。

例如，不存在阶为 $3675$ 的单群。若存在，令 $G$ 是阶为
\[
  3675=3\cdot5^2\cdot7^2
\]
的单群。$n_7$ 的唯一可能值为 $15$，故对 $Q\in\Syl_7(G)$，$|G:N_G(Q)|=15$，$|N_G(Q)|=245=5\cdot7^2$。令 $N=N_G(Q)$，$P\in\Syl_5(N)$。由 Sylow 同余条件应用于 $N$，得 $P\normal N$。由于 $|P|=5$，$P$ 不是 $G$ 的 Sylow $5$-子群，故包含在某个 $P^*\in\Syl_5(G)$ 中。由于 $P$ 是 $5$-群 $P^*$ 的指标 $5$ 子群，由定理 1，$P\normal P^*$，也就是说 $P^*\leq N_G(P)$。因此
\[
  \langle N,P^*\rangle\leq N_G(P),
  \qquad \text{从而 } 7^2\cdot5^2\mid |N_G(P)|.
\]
所以 $|G:N_G(P)|$ 整除 $3$。这不可能，因为 $P$ 不正规且 $G$ 没有指标 $3$ 的子群。

\subsection*{研究 Sylow $p$-子群交的正规化子}

前面第一种方法中的计数论证不能直接推广到非素数阶 Sylow 子群的原因之一是：若 $P\in\Syl_p(G)$ 且 $|P|=p^a$、$a\geq2$，则不同共轭的 $P$ 不一定只在单位子群处相交。若不同共轭的 $P$ 确实只交于单位子群，则仍可计数得到 $p$-幂阶元素个数为
\[
  n_p(|P|-1).
\]

假设存在 $R\in\Syl_p(G)$，$R\neq P$ 且 $P\cap R\neq 1$。令 $P_0=P\cap R$。则 $P_0<P$ 且 $P_0<R$，由定理 1，
\[
  P_0<N_P(P_0),
  \qquad
  P_0<N_R(P_0).
\]
可试图由此证明 $P_0$ 在 $G$ 中的正规化子足够大（即指标足够小），从而用前述方法得到矛盾。

一个特别有效的特殊情形是 $|P_0|=p^{a-1}$，也就是两个 Sylow $p$-子群 $R$ 与 $P$ 有很大的交。此时令 $N=N_G(P_0)$。由上述推理（因为 $P_0$ 是 $p$-群 $P$ 与 $R$ 的极大子群），$P_0\normal P$ 且 $P_0\normal R$，即
\[
  N\text{ 有两个不同的 Sylow }p\text{-子群 }P\text{ 与 }R.
\]
特别地，$|N|=p^a k$，其中由 Sylow 定理 $k\geq p+1$。

总之，若 Sylow $p$-子群两两交平凡，则可计数 $p$-幂阶元素；否则存在某个 Sylow 子群交，其正规化子“很大”。对任意群阶，事先未必能判断哪种现象发生，因此非单性证明可能需要分成两个互斥情形并分别推出矛盾。当 Sylow $p$-子群阶为 $p^2$ 时，这种过程尤其适用（例如，在证明阶为 $60$ 的单群同构于 $A_5$ 时，计数 $2$-幂阶元素就使用了这种思路——见第 4.5 节命题 23）。

下面给出一个引理，提供迫使非平凡 Sylow 交的充分条件。

\medskip
\noindent\textbf{引理 13.} 在有限群 $G$ 中，若 $n_p\not\equiv 1\pmod{p^2}$，则存在不同的 Sylow $p$-子群 $P,R$，使得 $P\cap R$ 在 $P$ 和 $R$ 中指标都为 $p$（因而在二者中正规）。

\begin{proof}
这是 Sylow 定理同余部分证明的一个简单细化。令 $P$ 通过共轭作用于集合 $\Syl_p(G)$。设 $O_1,
\ldots,O_s$ 为该作用的轨道，其中 $O_1=\{P\}$。若对每个 $R\in\Syl_p(G)$、$R\neq P$ 都有 $p^2\mid |P:P\cap R|$，则每个 $O_i$ 的大小均被 $p^2$ 整除（$i=2,
\ldots,s$）。于是作为轨道长度之和，$n_p=1+kp^2$，与假设矛盾。因此存在 $R\in\Syl_p(G)$ 使 $|P:P\cap R|=p$。
\end{proof}

例如，不存在阶为 $1053$ 的单群。若存在，令 $G$ 是阶为
\[
  1053=3^4\cdot13
\]
的单群，令 $P\in\Syl_3(G)$。必须有 $n_3=13$。但 $13\not\equiv1\pmod{3^2}$，所以存在 $P,R\in\Syl_3(G)$ 使 $|P\cap R|=3^3$。令 $N=N_G(P\cap R)$，则由上面的论证 $P,R\leq N$。因此 $3^4\mid |N|$ 且 $|N|>3^4$。唯一可能是 $N=G$，即 $P\cap R\normal G$，矛盾。

\subsection*{阶为 $168$ 的单群}

现在我们说明如何用许多技巧来剖析并分类某些单群。我们将分类阶为 $168$ 的单群。由于单群没有非平凡正规子群，这一过程不同于第 5.5 节的方法，但总体思想代表了有限单群研究中的常见方法。

假设存在阶为
\[
  168=2^3\cdot3\cdot7
\]
的单群 $G$。我们先推导它的许多性质：Sylow 子群的数目与结构、共轭类等。所有这些计算只依赖于 $G$ 的阶与单性。随后用这些结果证明 $G$ 唯一，最后证明阶为 $168$ 的单群确实存在。

因为 $|G|\nmid 6!$，有
\begin{enumerate}[label=(\arabic*)]
  \item $G$ 没有指标小于 $7$ 的真子群。
\end{enumerate}
否则 $G$ 作用于该子群的陪集会给出一个（因 $G$ 单而必为单射的）同态 $G\to S_n$，其中 $n\leq6$。

由 $G$ 的单性和 Sylow 定理立刻得到：
\begin{enumerate}[label=(\arabic*),resume]
  \item $n_7=8$，所以 Sylow $7$-子群的正规化子阶为 $21$。特别地，没有阶为 $2$ 的元素正规化 Sylow $7$-子群，且 $G$ 没有阶为 $14$ 的元素。
\end{enumerate}

若 $G$ 有阶为 $21$ 的元素，则 $G$ 的 Sylow $3$-子群正规化子的阶会被 $7$ 整除。于是 $n_3$ 与 $7$ 互素；又 $n_3\mid 8$，便有 $n_3=4$，与 (1) 矛盾。因此：
\begin{enumerate}[label=(\arabic*),resume]
  \item $G$ 没有阶为 $21$ 的元素。
\end{enumerate}

由 Sylow 定理，$n_3=7$ 或 $28$；下面排除前者。若 $n_3=7$，令 $P\in\Syl_3(G)$，并令 $T$ 是 $N_G(P)$（阶为 $24$）的 Sylow $2$-子群。每个 Sylow $3$-子群都正规化 $G$ 的某个 Sylow $7$-子群，因此 $P$ 正规化某个 $R\in\Syl_7(G)$。对每个 $t\in T$，$P=tPt^{-1}$ 也正规化 $tRt^{-1}$。子群 $T$ 通过共轭作用于八个 Sylow $7$-子群；由于由 (2) 没有阶为 $2$ 的元素正规化 Sylow $7$-子群，$T$ 的作用是传递的。于是 $P$ 正规化 $G$ 的每个 Sylow $7$-子群，即 $P$ 包含在所有 Sylow $7$-正规化子的交中。但这个交是 $G$ 的真正规子群，因此必须平凡，矛盾。故
\begin{enumerate}[label=(\arabic*),resume]
  \item $n_3=28$，Sylow $3$-子群的正规化子阶为 $6$。
\end{enumerate}

由于 $n_2=7$ 或 $21$，有 $n_2\not\equiv1\pmod 8$。由引理 13（或相应练习），存在一对不同 Sylow $2$-子群有非平凡交。选取所有这样的对中交的阶最大的 $T_1,T_2$，令 $U=T_1\cap T_2$。接下来证明
\begin{enumerate}[label=(\arabic*),resume]
  \item $U$ 是 Klein 四元群，且 $N_G(U)\cong S_4$。
\end{enumerate}

令 $N=N_G(U)$。由于 $|U|=2$ 或 $4$，且 $N$ 通过共轭置换 $U$ 的非单位元素，$N$ 中阶为 $7$ 的子群将与 $U$ 中某个阶为 $2$ 的元素交换，违背 (2)。故 $7\nmid |N|$。由引理 13 的思想，$N$ 有不止一个 Sylow $2$-子群，故 $|N|=2^a\cdot3$，其中 $a=2$ 或 $3$。令 $P\in\Syl_3(N)$。由于 $P$ 是 $G$ 的 Sylow $3$-子群，由 (4)，$N_N(P)$ 的阶为 $3$ 或 $6$。于是由 Sylow 定理，$N$ 必有四个 Sylow $3$-子群，并且 $N$ 通过共轭传递置换它们。由于任意阶为 $12$ 的群都有正规 Sylow $2$-子群或正规 Sylow $3$-子群（参见第 4.5 节），必有 $|N|=24$。令 $K$ 是 $N$ 作用于其四个 Sylow $3$-子群的共轭作用的核。若 $K=1$，则 $N\cong S_4$。若 $K\neq1$，因为 $K\leq N_N(P)$，$|K|$ 整除 $6$；且 $P$ 不正规化另一个 Sylow $3$-子群，所以 $P\nleq K$。于是 $|K|=2$。但 $N/K$ 是阶为 $12$ 的群，并且有不止一个 Sylow $2$-子群和四个 Sylow $3$-子群，矛盾于上面引用的阶 $12$ 群的性质。故 $N\cong S_4$。由于 $S_4$ 有唯一的非平凡正规 $2$-子群 $V_4$，得到 (5)。

由 $N\cong S_4$，$N$ 含有 $G$ 的 Sylow $2$-子群，而且 $N_N(P)\cong S_3$（由 (4) 也有 $N_G(P)\cong S_3$）。于是得到
\begin{enumerate}[label=(\arabic*),resume]
  \item $G$ 的 Sylow $2$-子群同构于 $D_8$；
  \item $G$ 中 Sylow $3$-子群的正规化子同构于 $S_3$，因此 $G$ 没有阶为 $6$ 的元素。
\end{enumerate}

由 (2) 和 (7)，没有阶为 $2$ 的元素与奇素数阶元素交换。若 $T\in\Syl_2(G)$，则由 (6)，$T\cong D_8$，所以 $Z(T)=\langle z\rangle$，其中 $z$ 阶为 $2$。于是 $T\leq C_G(z)$，而由刚才所述，$|C_G(z)|$ 没有奇素因子，故 $C_G(z)=T$。任意正规化 $T$ 的元素都正规化其中心，因而与 $z$ 交换，所以 Sylow $2$-子群自正规化。故
\begin{enumerate}[label=(\arabic*),resume]
  \item $n_2=21$，且若 $T\in\Syl_2(G)$、$Z(T)=\langle z\rangle$，则 $C_G(z)=T$。
\end{enumerate}

由于 $|C_G(z)|=8$，元素 $z$ 有 $21$ 个共轭。由 (6)，$G$ 有一个阶为 $4$ 的元素共轭类，其大小为 $42$。由 (2) 有 $48$ 个阶为 $7$ 的元素；由 (4) 有 $56$ 个阶为 $3$ 的元素。这些加起来已经占据了所有 $167$ 个非单位元素，因此每个阶为 $2$ 的元素都与 $z$ 共轭，即
\begin{enumerate}[label=(\arabic*),resume]
  \item $G$ 有唯一的阶为 $2$ 元素的共轭类。
\end{enumerate}

保持记号，令 $T\in\Syl_2(G)$，$U\leq T$，并令 $W$ 为 $T$ 中另一个 Klein 四元群。由 Sylow 定理可知 $U$ 与 $W$ 在 $G$ 中不共轭，因为它们在 $N_G(T)=T$ 中不共轭（参见第 4.5 节练习 50）。接下来证明
\begin{enumerate}[label=(\arabic*),resume]
  \item $N_G(W)\cong S_4$。
\end{enumerate}
为此写 $W=\langle z,w\rangle$，其中 $\langle z\rangle=Z(T)$。由于 $w$ 在 $G$ 中与 $z$ 共轭，$C_G(w)=T_0$ 是另一个含有 $W$ 且不同于 $T$ 的 Sylow $2$-子群。因此 $W=T\cap T_0$。因为 $U$ 是 Sylow $2$-子群最大交的任意选择，给出 (5) 的论证也给出 (10)。

现在记录已证明或易于推出的结果。

\medskip
\noindent\textbf{命题 14.} 若 $G$ 是阶为 $168$ 的单群，则：
\begin{enumerate}[label=(\arabic*)]
  \item $n_2=21$，$n_3=28$，$n_7=8$；
  \item $G$ 的 Sylow $2$-子群是二面体群，Sylow $3$-与 $7$-子群是循环群；
  \item $G$ 同构于 $A_7$ 的一个子群，并且没有指标 $\leq6$ 的子群；
  \item $G$ 的共轭类如下：单位元；两个阶为 $7$ 的元素类，各含 $24$ 个元素（可由任一阶为 $7$ 的元素及其逆元代表）；一个阶为 $3$ 的元素类，含 $56$ 个元素；一个阶为 $4$ 的元素类，含 $42$ 个元素；一个阶为 $2$ 的元素类，含 $21$ 个元素。特别地，$G$ 的每个元素阶都是某个素数的幂；
  \item 若 $T\in\Syl_2(G)$，而 $U,W$ 是 $T$ 中两个 Klein 四元群，则 $U,W$ 在 $G$ 中不共轭，且 $N_G(U)\cong N_G(W)\cong S_4$；
  \item $G$ 恰有三个极大子群的共轭类，其中两个同构于 $S_4$，另一个同构于阶为 $21$ 的非阿贝尔群。
\end{enumerate}

以上所有计算都建立在“存在阶为 $168$ 的单群”的假设上。没有导出矛盾并不证明这种群存在，只是强烈提示可能存在。下面说明 $G$ 的内部子群结构如何导出一个几何，使 $G$ 作用在其上，从而证明若阶为 $168$ 的单群存在则唯一，并最终证明它确实存在。

沿用记号，令 $U_1,
\ldots,U_7$ 为 $U$ 的共轭，$W_1,
\ldots,W_7$ 为 $W$ 的共轭。称 $U_i$ 为\emph{点}，$W_j$ 为\emph{直线}。定义一个“关联关系”：
\[
  \text{点 }U_i\text{ 在直线 }W_j\text{ 上}\quad\Longleftrightarrow\quad U_i\text{ 正规化 }W_j.
\]
注意 $U_i$ 正规化 $W_j$ 当且仅当 $U_iW_j\cong D_8$，而这又当且仅当 $W_j$ 正规化 $U_i$。在每个同构于 $S_4$ 的点或线稳定子中，有唯一正规 $4$-群 $V$，以及恰好三个其他（非正规）$4$-群 $A_1,A_2,A_3$。群 $VA_i$ 是该 $S_4$ 的三个 Sylow $2$-子群。因此
\begin{enumerate}[label=(\arabic*),resume]
  \item 每条直线恰含 $3$ 个点，每个点恰在 $3$ 条直线上。
\end{enumerate}
由于 $S_4$ 中任意两个非正规 $4$-群生成整个 $S_4$，从而唯一确定该 $S_4$ 中另外两个 Klein 四元群，得到
\begin{enumerate}[label=(\arabic*),resume]
  \item 直线上任意两个点唯一确定这条直线（也唯一确定其第三个点）。
\end{enumerate}
因为共有 $7$ 个点和 $7$ 条直线，简单计数给出
\begin{enumerate}[label=(\arabic*),resume]
  \item 任意两个点在唯一一条直线上，任意两条直线交于唯一一个点。
\end{enumerate}

这一点线配置因而满足所谓射影平面的公理。现在可直接证明该关联几何唯一，并可用图 \ref{fig:fano} 表示：点为顶点，直线为三角形的六条边和中线以及内切圆（见练习 27）。该关联几何称为\emph{阶为 $2$ 的射影平面}或\emph{Fano 平面}，记为 $\mathcal F$。一般地，一个“阶”为 $N$ 的射影平面有 $N^2+N+1$ 个点，并有同样数目的直线。注意在这一点上射影平面 $\mathcal F$ 的确存在——我们已经显式给出了满足 (11) 到 (13) 的点与直线——尽管群 $G$ 尚未知道存在。

\begin{figure}[htbp]
\centering
\begin{tikzpicture}[scale=2.0]
  \coordinate (A) at (90:1);
  \coordinate (B) at (210:1);
  \coordinate (C) at (330:1);
  \coordinate (D) at ($(B)!0.5!(C)$);
  \coordinate (E) at ($(A)!0.5!(C)$);
  \coordinate (F) at ($(A)!0.5!(B)$);
  \coordinate (O) at (0,0);
  \draw (A)--(B)--(C)--cycle;
  \draw (A)--(D);
  \draw (B)--(E);
  \draw (C)--(F);
  \draw (O) circle (0.5);
  \foreach \P in {A,B,C,D,E,F,O}{\fill (\P) circle (1.2pt);}
\end{tikzpicture}
\caption{图 1}
\label{fig:fano}
\end{figure}

该平面的一个自同构是点和直线的任意置换，要求保持关联关系。图中三角形的六个对称显然给出 $\mathcal F$ 的自同构，但 $\mathcal F$ 还有更多自同构。

每个 $g\in G$ 都通过共轭作用在点和线的集合上，且该作用保持关联关系。只有单位元固定所有点，因此通过此作用，$G$ 同构于 $\Aut(\mathcal F)$ 的一个子群。任何固定同一直线上的两个点以及不在该直线上的第三个点的 $\mathcal F$ 自同构都固定所有点。因此，$\mathcal F$ 的任一自同构由其在任意三个不共线点上的作用唯一确定。容易计算这种三元组有 $168$ 个，所以 $\mathcal F$ 至多有 $168$ 个自同构。这证明
\[
  \text{若单群 }G\text{ 存在，则它唯一，且 }G\cong\Aut(\mathcal F).
\]

分类过程还剩两步：证明 $\mathcal F$ 确有 $168$ 个自同构，并证明 $\Aut(\mathcal F)$ 确实是单群。虽然这两个步骤可以用图论方法完成，我们采用一种遵循“代数群”理论思想的方法。令 $V$ 是 $\FF_2$ 上三维向量空间，即 $V$ 是阶为 $8$ 的初等阿贝尔 $2$-群 $\ZZ_2\times\ZZ_2\times\ZZ_2$。由第 4.4 节命题 17，
\[
  \Aut(V)=\GL(V)\cong \GL_3(\FF_2)
\]
的阶为 $168$。称 $V$ 的七个 $1$ 维子空间（即非平凡循环子群）为点，称七个 $2$ 维子空间（即阶为 $4$ 的子群）为直线，并规定点 $p$ 与线 $L$ 关联当且仅当 $p\subseteq L$。这些点与线容易看出满足上面的 (11)--(13)，因而表示 Fano 平面。由于 $\GL(V)$ 忠实地作用在这些点和线上并保持关联，$|\Aut(\mathcal F)|\geq168$。结合已知上界，得到
\[
  \Aut(\mathcal F)\cong \GL(V)\cong\GL_3(\FF_2),
  \qquad |\Aut(\mathcal F)|=168.
\]

最后证明 $\GL(V)$ 是单群。反设 $H$ 是 $\GL(V)$ 的非平凡真正规子群。令 $\Omega$ 为七个点的集合，$N$ 为 $\GL(V)$ 中某个点的稳定子。由于 $\GL(V)$ 传递作用在 $\Omega$ 上，$N$ 指标为 $7$。所有 $N$ 的共轭的交固定所有点，因此是单位群。于是 $H\nleq N$，故 $\GL(V)=HN$。由 $|H:H\cap N|=|HN:N|$，得 $7\mid |H|$。由于 $\GL(V)$ 同构于 $S_7$ 的子群，且 $S_7$ 的 Sylow $7$-子群正规化子阶为 $42$，$\GL(V)$ 没有正规 Sylow $7$-子群，故由 Sylow 定理 $n_7(\GL(V))=8$。若 $H$ 有正规 Sylow $7$-子群，则它在 $H$ 中特征，因而在 $\GL(V)$ 中正规；所以 $H$ 也没有唯一 Sylow $7$-子群。由于 $n_7(H)\equiv1\pmod7$ 且 $n_7(H)\leq n_7(\GL(V))=8$，必有 $n_7(H)=8$。于是 $8\mid |H|$，从而 $56\mid |H|$；因 $H$ 真，必有 $|H|=56$。通常的计数论证说明 $H$ 有正规（故特征）Sylow $2$-子群，于是该子群在 $\GL(V)$ 中正规。但这会使 $\GL(V)$ 有唯一 Sylow $2$-子群。另一方面，$\GL_3(\FF_2)$ 中上三角矩阵组与下三角矩阵组都是阶为 $8$ 的子群，矛盾。于是证明如下定理。

\medskip
\noindent\textbf{定理 15.} 在同构意义下，阶为 $168$ 的单群唯一，即 $\GL_3(\FF_2)$；它也是射影平面 $\mathcal F$ 的自同构群。

注意，我们也可以把 $W_j$ 称为点、把 $U_i$ 称为直线。这种点线“对偶性”连同阶为 $168$ 单群的唯一性，可用来证明 $G$ 存在一个外自同构，它交换点与直线，即把 $U$ 共轭到 $W$。

许多有限单群族可用类似方法分类。在更一般的背景中，称为建筑的几何结构扮演射影平面的角色（射影平面是 $A_2$ 型建筑的特例）。在这种语境中，子群 $N_G(U)$ 与 $N_G(W)$ 是 $G$ 的抛物子群，而 $U,W$ 分别是它们的幺幂根基。特别地，所有单线性群都可由其抛物子群的结构与交来刻画，等价地，也可由它们在相关建筑上的作用来刻画。

\subsection*{关于群的存在性问题的注记}

如同数学其他领域（例如微分方程理论）中一样，我们可以假设某个数学系统存在（例如某方程的解），并由此推导关于该系统的大量信息。一般来说，如果经过相当努力仍未从初始假设推出矛盾，人们会开始怀疑确有系统满足这些假设。然而，任何数量的一致数据都不能证明存在性。假设我们分析一个阶为 $3^3\cdot7\cdot13\cdot409$ 的假想单群 $G$，类似于上面对阶为 $168$ 的单群的分析。经过一定工作可证明可能的 Sylow 数唯一：
\[
  n_3=7\cdot409,
  \qquad
  n_7=3^3\cdot13\cdot409,
  \qquad
  n_{13}=3^2\cdot7\cdot409,
  \qquad
  n_{409}=3^3\cdot7\cdot13.
\]
还可证明这样的 $G$ 没有阶为 $pq$（$p,q$ 为不同素数）的元素，没有阶为 $9$ 的元素，并且不同 Sylow 子群交平凡。于是计数所有 Sylow $p$-子群中的元素，恰好得到 $|G|$。此时我们会得到 $G$ 的完整子群结构与类方程，可能猜想存在这样的单群。但 Feit--Thompson 定理断言不存在奇合数阶单群。（不过要注意，阶为 $3^3\cdot7\cdot13\cdot409$ 的可能单群的这种配置正是 Feit--Thompson 定理证明中必须处理的情形之一，因此在这个例子中引用该结果实际上有循环之嫌。我们将在第 19.3 节证明不存在这种阶的单群；也见练习 29。）这里的要点是：即便我们拥有和阶 $168$ 情形一样多的数据，也不能在没有新技术的情况下证明存在。

当处理非单群时，我们至少有一种由较小群构造较大群的方法：半直积。尽管这种方法相当受限，它仍体现了非单群可以以某种构造方式由较小群拼合而成。对单群，这一过程完全失效；因此这条技术分界线强化了 H\"older 纲领的重要性：确定单群，并理解这些群如何组合成更大的群。

有限单群研究如阶 $168$ 的讨论所示，使用了许多与非单群研究相同的工具（用以剖析子群结构等），但还需要其他构造技术。正如上面提到的，这些技术常常涉及代数或几何方法，把单群构造成某些本身具有内在兴趣的数学结构的自同构群，从而以迷人的方式把群论与数学和科学的其他领域联系起来。因此，尽管我们在有限群分析中已经走了很远，但这一分支中仍有许多领域只是略作触及。

无限群的分析通常涉及相当不同的方法，下一节将介绍其中的一些。

\subsection*{练习}
\paragraph{计数元素}
\begin{enumerate}[label=\arabic*.]
  \item 证明：对固定 $P\in\Syl_p(G)$，若对所有 $R\in\Syl_p(G)-\{P\}$ 都有 $P\cap R=1$，则任意两个不同 Sylow $p$-子群 $P_1,P_2$ 都有 $P_1\cap P_2=1$。推出在这种情形下，$G$ 中非单位的 $p$-幂阶元素数为 $(|P|-1)|G:N_G(P)|$。
  \item 在群 $S_3\times S_3$ 中展示一对交为单位的 Sylow $2$-子群，并展示另一对交为阶 $2$ 的群的 Sylow $2$-子群。
  \item 证明若 $|G|=380$，则 $G$ 不单。[只需计数奇素数阶元素。]
  \item 证明不存在阶为 $80,351,3875,5313$ 的单群。
  \item 设 $G$ 是阶为 $pm$ 的可解群，其中 $p$ 是不整除 $m$ 的素数，$P\in\Syl_p(G)$。若 $N_G(P)=P$，证明 $G$ 有一个阶为 $m$ 的正规子群。证明中哪里用到了 $G$ 的可解性？
\end{enumerate}

\paragraph{利用小指标子群}
\begin{enumerate}[label=\arabic*.,resume]
  \item 证明不存在阶为 $2205,4125,5103,6545,6435$ 的单群。
\end{enumerate}

\paragraph{置换表示}
\begin{enumerate}[label=\arabic*.,resume]
  \item 证明不存在阶为 $1755$ 或 $5265$ 的单群。[用 Sylow $3$-子群说明 $G\leq S_{13}$，并考察 Sylow $13$-子群的正规化子。]
  \item 证明不存在阶为 $792$ 或 $918$ 的单群。
  \item 证明不存在阶为 $336$ 的单群。
\end{enumerate}

\paragraph{让不同 $p$-子群相互制约}
\begin{enumerate}[label=\arabic*.,resume]
  \item 证明不存在阶为 $4095,4389,5313,6669$ 的单群。
  \item 证明不存在阶为 $4851$ 或 $5145$ 的单群。
  \item 证明不存在阶为 $9555$ 的单群。[令 $Q\in\Syl_{13}(G)$，$P\in\Syl_7(N_G(Q))$。证明 $Q\normal N_G(P)$，并说明为何矛盾。]
\end{enumerate}

\paragraph{Sylow 交的正规化子}
\begin{enumerate}[label=\arabic*.,resume]
  \item 设 $G$ 有不止一个 Sylow $p$-子群。取一对不同 Sylow $p$-子群 $P,Q$，使 $|P\cap Q|$ 在所有这样的对中最大。证明 $N_G(P\cap Q)$ 有不止一个 Sylow $p$-子群，且 $N_G(P\cap Q)$ 的任意两个不同 Sylow $p$-子群都交于 $P\cap Q$。（因此 $|N_G(P\cap Q)|$ 被 $p\cdot |P\cap Q|$ 以及某个不同于 $p$ 的素数整除。注意 $N_G(P\cap Q)$ 的 Sylow $p$-子群不一定是 $G$ 的 Sylow 子群。）
  \item 证明不存在阶为 $144,525,2025,3159$ 的单群。
\end{enumerate}

\paragraph{一般练习}
\begin{enumerate}[label=\arabic*.,resume]
  \item 分类阶为 $105$ 的群。
  \item 证明不存在奇阶 $<10000$ 的非阿贝尔单群。
  \item \begin{enumerate}[label=(\alph*)]
    \item 证明不存在阶为 $420$ 的单群。
    \item 证明除阶 $2,60,168,360$ 外，不存在偶阶 $<500$ 的单群。
  \end{enumerate}
  \item 证明若 $G$ 阶为 $36$，则 $G$ 有正规 Sylow $2$-子群或正规 Sylow $3$-子群。
  \item 证明没有阶为 $6$ 子群的阶 $12$ 群同构于 $A_4$。
  \item 证明没有阶为 $6$ 元素的阶 $24$ 群同构于 $S_4$。
  \item 推广引理 13：证明若 $n_p\not\equiv1\pmod{p^k}$，则存在不同 Sylow $p$-子群 $P,R$，使 $P\cap R$ 在二者中指标至多为 $p^{k-1}$。
  \item 在所有不同 Sylow $p$-子群对中，取 $P,R$ 使 $|P\cap R|$ 最大。证明 $N_G(P\cap R)$ 不是 $p$-群。
  \item 设 $A,B$ 是 $G$ 的 Sylow $p$-子群 $P$ 的正规子集。证明若 $A,B$ 在 $G$ 中共轭，则它们在 $N_G(P)$ 中共轭。
  \item 设 $G$ 是阶为 $pqr$ 的群，其中 $p<q<r$ 为素数。证明 $G$ 的 Sylow $r$-子群正规。
  \item 设 $G$ 是阶为 $p^2qr$ 的单群，其中 $p,q,r$ 为素数。证明 $|G|=60$。
  \item 判断并证明或给出反例：若 $G$ 是阶为 $168$ 的群且有不止一个 Sylow $7$-子群，则 $G$ 单。
  \item 证明若 $\mathcal F$ 是任意满足阶 $168$ 单群小节中 (11)--(13) 的点线集合，则 $\mathcal F$ 的关联图唯一确定，并与图 \ref{fig:fano} 相同（至多重标记点和线）。[取一条直线和一个不在该直线上的点。把这条直线画成一个等边三角形的底边，把该点画成不在底边上的顶点。用公理证明剩余点与直线的关联关系于是如图 \ref{fig:fano} 那样唯一确定。]
  \item 设 $G$ 是阶为 $3^3\cdot7\cdot13\cdot409$ 的单群。对每个 $p\in\{3,7,13,409\}$ 计算 $n_p$ 的所有允许值，并化简到每个 $n_p$ 都只有唯一可能值的情形。
  \item 利用假想阶为 $3^3\cdot7\cdot13\cdot409$ 单群的 Sylow 数信息，证明不存在这样的群。[使用次数为 $819$ 的置换表示。]
  \item 假设 $G$ 是阶为 $720$ 的单群。尽可能多地求出 $G$ 的性质（Sylow 数、Sylow 子群同构类型、共轭类等）。是否存在这样的群？
\end{enumerate}

\section{自由群简介}

本节介绍所谓自由群的基本理论。这将使我们能够精确化前几章中使用过的“生成元”和“关系”的概念。本节结果只依赖同态的基本理论。

由集合 $S$ 生成的自由群 $F(S)$ 的基本思想是：$S$ 中元素不满足任何关系（$S$ 是“无关系”的）。例如，若 $S=\{a,b\}$，则两个生成元 $a,b$ 上的自由群的元素形如 $a,aa,ab,abab,ba b$ 等，称为 $a,b$ 中的字，还包括这些元素的逆；所有这些元素都视为不同。若把相同项合并，可得到熟悉形式如 $a^3b^{-2},aba^{-1}b^2$ 等。这样的元素通过连接它们的字相乘。直观上可先把 $F(S)$ 定义为 $S$ 中所有字的集合，并以连接作为乘法。虽然本质上我们就是这样做，但为了证明连接运算良定义且满足结合律，需要更形式化的构造。毕竟，即使熟悉的记号 $a^n$ 表示 $a\cdot a\cdots a$（$n$ 项），也只有在知道乘积与括号放置无关后才合法。

自由群没有关系这一事实的一个重要性质是：从集合 $S$ 到任意群 $G$ 的任意映射，都唯一地扩张为从群 $F(S)$ 到 $G$ 的同态。基本原因是我们指定了生成元的像，而其他元素的像由同态性质唯一决定；没有关系需要担心，因而生成元的像可任意指定。这通常称为自由群的泛性质，并且事实上刻画了 $F(S)$。

“自由性”的概念出现在许多代数系统中。在线性代数中，当代数系统是向量空间时，$F(S)$ 就是以 $S$ 为基的向量空间。每个向量都是 $S$ 中元素的唯一线性组合（类似于“字”），而从基 $S$ 到另一个向量空间 $V$ 的任意集合映射都唯一扩张为线性变换。

在交换图语言中，自由群的泛性质可写作：给定集合映射 $\varphi:S\to G$，其中 $G$ 是群，存在唯一同态 $\Phi:F(S)\to G$ 使 $\Phi|_S=\varphi$，即图
\[
\begin{array}{ccc}
S & \longrightarrow & F(S)\\
 & \searrow_{\varphi} & \downarrow^{\Phi}\\
 & & G
\end{array}
\]
交换。

下面构造 $F(S)$。构造中的唯一困难是验证字的连接运算良定义且满足结合律。为证明结合律，我们回到最基本层面：字中所有指数都是 $\pm1$。

先为 $S$ 的元素引入逆和单位元。令 $S^{-1}$ 为一个与 $S$ 不交的集合，并有一个从 $S$ 到 $S^{-1}$ 的双射。对每个 $s\in S$，把相应元素记为 $s^{-1}$；对每个 $t\in S^{-1}$，把相应的 $S$ 中元素记为 $t^{-1}$，于是 $(s^{-1})^{-1}=s$。再取一个不在 $S\cup S^{-1}$ 中的单点集 $\{1\}$，令 $1^{-1}=1$，并对任意 $x\in S\cup S^{-1}\cup\{1\}$ 令 $(x^{-1})^{-1}=x$。

$S$ 上的一个\emph{字}定义为一个序列
\[
  (s_1,s_2,s_3,\ldots),
  \qquad s_i\in S\cup S^{-1}\cup\{1\},
\]
且对充分大的 $i$ 有 $s_i=1$。换言之，可把一个字看作由 $S$ 中元素及其逆组成的有限乘积（允许重复）。为保证表达唯一，我们只考虑相邻项之间没有显然消去的字。字 $(s_1,s_2,s_3,\ldots)$ 称为\emph{约化的}，若
\begin{enumerate}[label=(\arabic*)]
  \item 当 $s_i\neq1$ 时，$s_{i+1}\neq s_i^{-1}$；
  \item 若某个 $k$ 有 $s_k=1$，则对所有 $i\geq k$，$s_i=1$。
\end{enumerate}
约化字 $(1,1,1,\ldots)$ 称为空字，记为 $1$。以后把约化字
\[
  (s_1^{\varepsilon_1},s_2^{\varepsilon_2},\ldots,s_n^{\varepsilon_n},1,1,\ldots)
\]
简写为
\[
  s_1^{\varepsilon_1}s_2^{\varepsilon_2}\cdots s_n^{\varepsilon_n},
  \qquad s_i\in S,
  \quad \varepsilon_i=\pm1.
\]
令 $F(S)$ 为 $S$ 上所有约化字的集合，并通过
\[
  s\longmapsto (s,1,1,1,\ldots)
\]
把 $S$ 嵌入 $F(S)$。在此嵌入下把 $S$ 与其像等同。若 $S=\varnothing$，则 $F(S)=\{1\}$。

现在在 $F(S)$ 上定义二元运算。技术困难是保证两个约化字的乘积仍为约化字。定义虽然形式化，但只是“相邻逆元逐次消去”的规则（例如，$ab^{-1}a$ 乘以 $a^{-1}ba$ 应约化为 $aa$）。设
\[
  r_1^{\delta_1}\cdots r_m^{\delta_m},
  \qquad
  s_1^{\varepsilon_1}\cdots s_n^{\varepsilon_n}
\]
为约化字，并先假设 $m\leq n$。令 $k$ 为范围 $1\leq k\leq m+1$ 中使
\[
  s_k^{\varepsilon_k}\neq r_{m-k+1}^{-\delta_{m-k+1}}
\]
成立的最小整数。则这两个约化字的乘积定义为
\[
\begin{cases}
 r_1^{\delta_1}\cdots r_{m-k+1}^{\delta_{m-k+1}}
 s_k^{\varepsilon_k}s_{k+1}^{\varepsilon_{k+1}}\cdots s_n^{\varepsilon_n}, & k\leq m,\\[2mm]
 s_{m+1}^{\varepsilon_{m+1}}\cdots s_n^{\varepsilon_n}, & k=m+1\leq n,\\[2mm]
 1, & k=m+1\text{ 且 }m=n.
\end{cases}
\]
当 $m\geq n$ 时类似地定义乘积。因此无论哪种情形，乘积都是一个约化字，由此得到 $F(S)$ 上的二元运算。

\medskip
\noindent\textbf{定理 16.} $F(S)$ 在上述二元运算下是一个群。

\begin{proof}
容易检查空字 $1$ 是单位元，约化字
\[
  s_1^{\varepsilon_1}s_2^{\varepsilon_2}\cdots s_n^{\varepsilon_n}
\]
的逆为
\[
  s_n^{-\varepsilon_n}s_{n-1}^{-\varepsilon_{n-1}}\cdots s_1^{-\varepsilon_1}.
\]
证明的困难部分是结合律。可对所涉字的“长度”作归纳并分情况讨论；也可如下进行。对每个 $s\in S\cup S^{-1}\cup\{1\}$，定义 $\alpha_s:F(S)\to F(S)$ 为
\[
\alpha_s(s_1^{\varepsilon_1}\cdots s_n^{\varepsilon_n})=
\begin{cases}
 ss_1^{\varepsilon_1}\cdots s_n^{\varepsilon_n}, & s_1^{\varepsilon_1}\neq s^{-1},\\[1mm]
 s_2^{\varepsilon_2}\cdots s_n^{\varepsilon_n}, & s_1^{\varepsilon_1}=s^{-1}.
\end{cases}
\]
由于 $\alpha_{s^{-1}}\circ\alpha_s$ 是 $F(S)$ 到自身的恒等映射，$\alpha_s$ 是 $F(S)$ 的一个置换。令 $A(S)$ 为集合 $F(S)$ 上的对称群中由 $\{\alpha_s\mid s\in S\}$ 生成的子群。容易看出映射
\[
  s_1^{\varepsilon_1}s_2^{\varepsilon_2}\cdots s_n^{\varepsilon_n}
  \longmapsto
  \alpha_{s_1^{\varepsilon_1}}\circ\alpha_{s_2^{\varepsilon_2}}\circ\cdots\circ\alpha_{s_n^{\varepsilon_n}}
\]
是 $F(S)$ 与 $A(S)$ 之间的集合双射，并保持它们的二元运算。由于 $A(S)$ 是群，因而满足结合律，所以 $F(S)$ 也满足结合律。
\end{proof}

自由群的泛性质现在立即得到。

\medskip
\noindent\textbf{定理 17.} 设 $G$ 是群，$S$ 是集合，$\varphi:S\to G$ 是集合映射。则存在唯一的群同态 $\Phi:F(S)\to G$ 使上述图交换。

\begin{proof}
若 $\Phi$ 是同态，则必须满足
\[
  \Phi(s_1^{\varepsilon_1}s_2^{\varepsilon_2}\cdots s_n^{\varepsilon_n})
  =\varphi(s_1)^{\varepsilon_1}\varphi(s_2)^{\varepsilon_2}\cdots\varphi(s_n)^{\varepsilon_n},
\]
这证明唯一性。直接检查该公式确实定义了一个同态，得到存在性。
\end{proof}

\medskip
\noindent\textbf{推论 18.} $F(S)$ 在保持集合 $S$ 上恒等的唯一同构意义下唯一。

\begin{proof}
这由泛性质推出。设 $F(S)$ 与 $F'(S)$ 是两个由 $S$ 生成的自由群。由于 $S$ 包含于二者，由定理的泛性质得到唯一同态
\[
  \Phi:F(S)\to F'(S),
  \qquad
  \Phi':F'(S)\to F(S),
\]
且二者在 $S$ 上都是恒等。复合 $\Phi'\Phi$ 是 $F(S)$ 到自身、在 $S$ 上恒等的同态；由唯一性，它必须是恒等映射。同理 $\Phi\Phi'$ 也是恒等。因此 $\Phi$ 是同构。
\end{proof}

\medskip
\noindent\textbf{定义.} 群 $F(S)$ 称为集合 $S$ 上的\emph{自由群}。若存在集合 $S$ 使 $F=F(S)$，则称 $F$ 是\emph{自由群}；此时称 $S$ 为 $F$ 的一组\emph{自由生成元}（或\emph{自由基}）。$S$ 的基数称为该自由群的\emph{秩}。

现在可以用指数记号简化自由群中的表达式，例如写 $a^3b^{-2}$ 代替形式约化字 $aaab^{-1}b^{-1}$。不过，在 $\{a,b\}$ 上的自由群中，像 $aba$ 这样的表达式不能再简化。下面提到该领域中的一个重要定理。

\medskip
\noindent\textbf{定理 19（Schreier）.} 自由群的子群仍是自由群。

这个定理的证明并不平凡，这里不包括证明。使用覆盖空间可以给出一个漂亮的证明（参见 J.-P. Serre, \emph{Trees}, Springer-Verlag, 1980）。

\subsection*{表示}

设 $G$ 是任意群。则 $G$ 是某个自由群的同态像：取 $S=G$，并取 $\varphi$ 为从集合 $G$ 到群 $G$ 的恒等映射；定理 17 给出一个从 $F(G)$ 到 $G$ 的满同态。更一般地，若 $S$ 是 $G$ 的任意生成子集，即 $G=\langle S\rangle$，则也存在唯一的满同态 $F(S)\to G$，它在 $S$ 上为恒等。注意现在可独立表述“子集生成群”：$G=\langle S\rangle$ 当且仅当把 $S$ 到 $G$ 的恒等映射扩张而得的映射 $\pi:F(S)\to G$ 是满射。

\medskip
\noindent\textbf{定义.} 设 $S$ 是群 $G$ 的子集，且 $G=\langle S\rangle$。
\begin{enumerate}[label=(\arabic*)]
  \item $G$ 的一个\emph{表示}是一个二元组 $(S,R)$，其中 $R$ 是 $F(S)$ 中的一组字，使得 $F(S)$ 中 $\langle R\rangle$ 的正规闭包（即包含 $\langle R\rangle$ 的最小正规子群）等于同态 $\pi:F(S)\to G$ 的核（这里 $\pi$ 扩张 $S$ 到 $S$ 的恒等映射）。$S$ 的元素称为 $G$ 的\emph{生成元}，$R$ 的元素称为 $G$ 的\emph{关系}。
  \item 若存在一个表示 $(S,R)$ 使 $S$ 有限，则称 $G$ 是\emph{有限生成的}；若存在一个表示 $(S,R)$ 使 $S$ 与 $R$ 都有限，则称 $G$ 是\emph{有限表示的}。
\end{enumerate}

注意，若 $(S,R)$ 是一个表示，则映射 $F(S)\to G$ 的核不是 $\langle R\rangle$ 本身，而是由 $R$ 及其所有共轭生成的更大的正规闭包。即使固定集合 $S$，一个群也会有许多不同表示（例如可以加入冗余关系）。若 $G$ 有有限表示，$S=\{s_1,
\ldots,s_n\}$ 且 $R=\{w_1,
\ldots,w_k\}$，则按前几章写作
\[
  G=\langle s_1,s_2,\ldots,s_n\mid w_1=w_2=\cdots=w_k=1\rangle.
\]
若 $w$ 是字 $w_1w_2^{-1}$，则把 $w=1$ 写作 $w_1=w_2$。

\paragraph{例}
\begin{enumerate}[label=(\arabic*)]
  \item 每个有限群都是有限表示的。设 $G=\{g_1,
\ldots,g_n\}$。取 $S=G$，令 $\pi:F(S)\to G$ 扩张 $S$ 上的恒等映射。令 $R_0$ 为所有字 $g_ig_jg_k^{-1}$ 的集合，其中 $g_ig_j=g_k$ 在 $G$ 中成立。显然 $R_0\subseteq\ker\pi$。若 $N$ 是 $R_0$ 在 $F(S)$ 中的正规闭包，并令 $\overline G=F(S)/N$，则 $G$ 是 $\overline G$ 的同态像。而且元素集合 $\{\overline g_i\mid i=1,
\ldots,n\}$ 在乘法下封闭；由于它生成 $\overline G$，必有 $|\overline G|\leq |G|$，从而 $N=\ker\pi$，所以 $(S,R_0)$ 是 $G$ 的一个表示。

  这说明给定有限群 $G$ 的一个表示 $(S,R)$ 的充分条件：
  \begin{enumerate}[label=(\roman*)]
    \item $S$ 必须生成 $G$；
    \item 任意由 $S$ 生成并满足 $R$ 中关系的群，其阶都不超过 $|G|$。
  \end{enumerate}

  \item 阿贝尔群很容易表示。例如
  \[
    \ZZ\cong F(\{a\})=\langle a\rangle,
  \]
  \[
    \ZZ\times\ZZ\cong \langle a,b\mid [a,b]=1\rangle,
  \]
  \[
    \ZZ_n\times\ZZ_m\cong \langle a,b\mid a^n=b^m=[a,b]=1\rangle.
  \]
  回忆 $[a,b]=a^{-1}b^{-1}ab$。

  \item 前几章中一些熟悉的非阿贝尔群有简单表示：
  \[
    D_{2n}=\langle r,s\mid r^n=s^2=1,\ s^{-1}rs=r^{-1}\rangle,
  \]
  \[
    Q_8=\langle i,j\mid i^4=1,
\ j^2=i^2,
\ j^{-1}ij=i^{-1}\rangle.
  \]
  例如，要检查 $D_{2n}$ 的表示，注意这些关系说明该群有一个由 $r$ 生成、阶至多为 $n$ 的正规子群，其商由 $s$ 生成且阶至多为 $2$。因此任何满足这些生成元与关系的群阶至多为 $2n$。由于我们已经知道存在阶为 $2n$ 并满足这些条件的群 $D_{2n}$，故该抽象表示就是 $D_{2n}$。

  \item 如第 1.2 节所述，一般而言，甚至判断给定生成元与关系所定义的群是否为单位群都极其困难。例如下面两个表示中，第一个群是无限群，第二个是单位群：
  \[
  \langle x_1,x_2,x_3,x_4\mid x_2^{-1}x_1x_2=x_1^2,
  x_3^{-1}x_2x_3=x_2^2,
  x_4^{-1}x_3x_4=x_3^2,
  x_1^{-1}x_4x_1=x_4^2\rangle,
  \]
  \[
  \langle x_1,x_2,x_3\mid x_2x_1x_2^{-1}=x_1^2,
  x_3x_2x_3^{-1}=x_2^2,
  x_1x_3x_1^{-1}=x_3^2\rangle.
  \]

  \item 容易看出 $S_n$ 由相邻换位 $(12),(23),\ldots,(n-1\ n)$ 生成，并且满足关系
  \[
    ((i\ i+1)(i+1\ i+2))^3=1,
  \]
  以及
  \[
    [(i\ i+1),(j\ j+1)]=1\quad\text{当 }|i-j|\geq2.
  \]
  可对 $n$ 归纳证明这些关系构成 $S_n$ 的一个表示：
  \[
    S_n\cong\langle t_1,
\ldots,t_{n-1}\mid t_i^2=1,
\ (t_it_{i+1})^3=1,
\ [t_i,t_j]=1\text{ whenever }|i-j|\geq2\rangle.
  \]
\end{enumerate}

如第 1.6 节所述，可用群的表示寻找群间同态或群的自同构。例如在分类阶为 $6$ 的群时，我们证明任意非阿贝尔阶 $6$ 群由一个阶为 $3$ 的元素和一个阶为 $2$、使其取逆的元素生成；因此存在从 $S_3$ 到任意非阿贝尔阶 $6$ 群的满同态（比较阶后为同构）。更一般地，设 $G$ 由生成元 $a,b$ 和关系 $r_1,
\ldots,r_k$ 表示。若 $H$ 中元素 $a',b'$ 满足这些关系，则存在从 $G$ 到 $H$ 的同态。具体地，若 $\pi:F(\{a,b\})\to G$ 是表示同态，则可定义 $\pi':F(\{a,b\})\to H$，使 $\pi'(a)=a'$ 且 $\pi'(b)=b'$。由于 $a',b'$ 满足 $G$ 的关系，有 $\ker\pi\leq\ker\pi'$，所以 $\pi'$ 通过 $\ker\pi$ 分解，从而得到
\[
  G\cong F(\{a,b\})/\ker\pi\longrightarrow H.
\]
特别地，若 $\langle a',b'\rangle=H=G$，则该同态是 $G$ 的自同构。反过来，任一自同构必须把一组生成元送到另一组满足相同关系的生成元。

例如
\[
  D_8=\langle a,b\mid a^2=b^4=1,
\ aba=b^{-1}\rangle.
\]
任意一对元素 $a',b'$，其中 $a'$ 是非中心的阶 $2$ 元素，$b'$ 阶为 $4$，都满足相同关系。$D_8$ 有四个非中心阶 $2$ 元素和两个阶 $4$ 元素，因此 $D_8$ 有 $8$ 个自同构。

同样，$Q_8$ 中任意一对阶为 $4$、且不相等也不互为逆的元素必生成 $Q_8$，并满足例 3 中的关系。容易检查这样的对有 $24$ 个，所以
\[
  |\Aut(Q_8)|=24.
\]

自由对象也可在许多（但不是所有）其他范畴中构造。例如，幺半群是带有二元运算、满足除逆元存在性外所有群公理的集合。幺半群范畴中的自由对象在理论计算机科学中起基础作用，用以刻画机器（Turing 机等）的行为。后面章节还会遇到自由代数（即多项式代数）与自由模。

\subsection*{练习}
\begin{enumerate}[label=\arabic*.]
  \item 设 $F_1,F_2$ 是有限秩自由群。证明 $F_1\cong F_2$ 当且仅当它们有相同秩。若要把证明推广到无限秩（结论仍为真），还需要哪些事实？
  \item 证明若 $|S|>1$，则 $F(S)$ 非阿贝尔。
  \item 证明两个生成元上的自由群的换位子子群不是有限生成的（特别地，有限生成群的子群不必有限生成）。
  \item 证明自由群的每个非单位元素都有无限阶。
  \item 用两个生成元给出 $A_4$ 的一个有限表示。
  \item 用两个生成元给出 $S_4$ 的一个有限表示。
  \item 证明以下是阶为 $8$ 的四元数群的一个表示：
  \[
    Q_8=\langle a,b\mid a^2=b^2,
\ a^{-1}ba=b^{-1}\rangle.
  \]
  \item 用表示求群 $\ZZ_2\times\ZZ_4$ 和 $\ZZ_4\times\ZZ_4$ 的自同构群的阶。
  \item 证明 $\Aut(Q_8)\cong S_4$。
  \item 本题展示 $S_6$ 的一个非内自同构（因此，连同第 4.4 节练习 19，说明 $|\Aut(S_6):\Inn(S_6)|=2$）。令
  \[
  \begin{aligned}
  t_1'&=(12)(34)(56),&
  t_2'&=(14)(25)(36),\\
  t_3'&=(13)(24)(56),&
  t_4'&=(12)(36)(45),\\
  t_5'&=(14)(23)(56).
  \end{aligned}
  \]
  证明 $t_1',
\ldots,t_5'$ 满足：$(t_i')^2=1$；当 $|i-j|\geq2$ 时 $(t_i't_j')^2=1$；以及对 $i=1,2,3,4$，$(t_i't_{i+1}')^3=1$。推出映射
  \[
    (12)\mapsto t_1',\quad (23)\mapsto t_2',\quad (34)\mapsto t_3',\quad
    (45)\mapsto t_4',\quad (56)\mapsto t_5'
  \]
  扩张为 $S_6$ 的一个自同构，且它显然不是内自同构，因为它不把换位送到换位。
  \item 设 $S$ 是集合。表示 $\langle S\mid [s,t]=1\text{ for all }s,t\in S\rangle$ 所给的群称为 $S$ 上的自由阿贝尔群，记为 $A(S)$。证明 $A(S)$ 有如下泛性质：若 $G$ 是任意阿贝尔群，$\varphi:S\to G$ 是任意集合映射，则存在唯一群同态 $\Phi:A(S)\to G$ 使 $\Phi|_S=\varphi$。推出若 $A$ 是基数为 $n$ 的集合上的自由阿贝尔群，则
  \[
    A\cong \underbrace{\ZZ\times\ZZ\times\cdots\times\ZZ}_{n\text{ 个因子}}.
  \]
  \item 设 $S$ 是集合，$c$ 是正整数。表述 $S$ 上幂零类为 $c$ 的自由幂零群的概念，并证明它相对于幂零类 $\leq c$ 的幂零群具有适当的泛性质。
  \item 证明不存在一个由两个元素生成的幂零群 $N$，使得每个由两个元素生成的幂零群都是 $N$ 的同态像（也就是说，前一题中指定类 $c$ 是必要的）。
\end{enumerate}
